\documentclass[12pt,openany]{book}

\usepackage[margin=0.9in]{geometry}
\usepackage{amsmath,amssymb,amsthm}
\usepackage{graphicx}
\usepackage[T1]{fontenc}
\usepackage[hidelinks]{hyperref}
\usepackage[spanish]{babel}

\setlength{\parindent}{0pt}
\setlength{\parskip}{4pt}

\theoremstyle{definition}
\newtheorem{definition}{Definición}[chapter]

\usepackage[hidelinks]{hyperref}
\hypersetup{
  pdftitle={Transformaciones Lineales, Geométricas y Estructurales: Una Perspectiva Axiomática y Libre de Coordenadas},
  pdfauthor={Jos\'e Portillo},
  pdfsubject={Álgebra lineal, geometría, tensores y estructura logarítmica},
  pdfkeywords={transformaciones lineales, productos internos, tensores, producto tensorial, logaritmos}
}

\begin{document}

\pagestyle{plain}

\frontmatter

\begin{titlepage}
\thispagestyle{empty}

\vspace*{2.2cm}

{\Large\bfseries
Transformaciones Lineales, Geométricas y Estructurales\par}
\vspace{0.4cm}

\vspace{1.4cm}
{\Large Jos\'e Portillo\par}

\vspace{1.8cm}
\rule{\textwidth}{0.6pt}
\vspace{0.35cm}

{\small
\textbf{Fecha de publicación:} 15 de febrero de 2026\par
\textbf{Versión:} 0.4\par
}

\vfill

\end{titlepage}


\thispagestyle{empty}
\null\vfill
\begin{center}
A la memoria de\par
Mercelena Portillo Holskin,\par
Eldon C. Hall,\par
Hugh Blair-Smith\par
\vspace{1cm}
A mi familia
\end{center}
\vfill\null
\clearpage

\tableofcontents

\chapter*{Prefacio}
\addcontentsline{toc}{chapter}{Prefacio}

Todo algoritmo, por sofisticado que sea, se reduce a una
secuencia finita de operaciones elementales: sumas, multiplicaciones,
comparaciones.
Todo sistema físico que realiza cómputo --- ya sea construido con válvulas de
vacío, transistores, guías de onda ópticas o cualquier sustrato aún por
concebir --- implementa esas mismas operaciones a través de medios físicos
distintos.
Lo que persiste (subyace) a través de todos esos cambios de medio y de representación es
una capa de estructura matemática.
Esta monografía se ocupa de esa capa que permanece invariante.

\medskip

La pregunta que organiza la exposición no es, por tanto, \emph{cómo} se lleva
a cabo un cómputo dado, sino \emph{qué} se está computando --- y qué
propiedades estructurales posee ese objeto con independencia del sistema de
coordenadas, la base o el hardware en que acontezca estar expresado.
Esta monografía nació como un intento de comprender qué permanece invariante
cuando se eliminan las fórmulas, las coordenadas y las representaciones.

\medskip

A lo largo del texto se establecerá de forma recurrente una distinción entre
una teoría axiomática y un modelo concreto de esa teoría.
En un tratamiento axiomático riguroso de la geometría euclidiana, objetos como
\textit{punto}, \textit{recta} y \textit{plano} se toman como términos
primitivos (no definidos), y los axiomas especifican únicamente las relaciones
que deben satisfacer.
Un sistema de coordenadas cartesianas no redefine estos primitivos; más bien,
proporciona un \emph{modelo} en el que los puntos se representan como tuplas
en~$\mathbb{R}^n$, las rectas como conjuntos solución de ecuaciones lineales,
y la incidencia se convierte en pertenencia a un conjunto~\cite{iribarren2}.
Como se afirma en~\cite{wilder}, esto es \textit{la asignación de
significados}.
Esta perspectiva sitúa las fórmulas coordenadas como representaciones de la
estructura, y no como la estructura misma.
La misma distinción atraviesa todos los capítulos de este texto.

\medskip

Los espacios vectoriales se introducen como conjuntos abstractos dotados de
adición y multiplicación por escalar.
Las transformaciones lineales se definen como las aplicaciones que preservan
dicha estructura.

\medskip

Al introducir formas bilineales y, en última instancia, productos internos, se
permite que los vectores interactúen de un modo que produce escalares.
De esta única interacción emergen normas, ortogonalidad, distancias y objetos
geométricos clásicos como bolas, esferas y elipsoides.

\medskip

Muchas de las construcciones encontradas anteriormente --- formas bilineales,
productos internos, multiplicación de matrices --- comparten un rasgo común:
son lineales en cada argumento por separado.
Esta observación conduce de forma natural a las aplicaciones multilineales y a
los tensores.

\medskip

Los tensores se tratan como objetos multilineales abstractos definidos sin
coordenadas.
El producto tensorial se introduce mediante su propiedad universal, que explica
cómo los problemas multilineales se reducen canónicamente a problemas lineales.
La notación indicial, los arreglos numéricos y las implementaciones
computacionales se presentan explícitamente como representaciones.

\medskip

El texto se organiza en torno a cuatro temas estructurales sucesivos:
estructura lineal, estructura geométrica y topológica, estructura multilineal
y estructura logarítmica.
Cada tema es un refinamiento de la misma idea organizadora: la estructura se
revela a través de la invariancia bajo transformaciones.
El tema logarítmico ilustra, en particular, que este principio no está
confinado al ámbito lineal --- el logaritmo es un isomorfismo global entre
estructuras multiplicativas y aditivas, y como tal pertenece al mismo marco
invariante que los vectores y los tensores.

\medskip

Se asume que el lector posee familiaridad previa con teoría de conjuntos
básica~\cite{wilder}\cite{cohen}\cite{apostol}, relaciones, funciones
\cite{apostol}, $n$-tuplas~\cite{iribarren}\cite{cohen} y álgebra
elemental~\cite{iribarren}.
Se asume asimismo conocimiento básico de cálculo~\cite{apostol}.
Los reportes~\cite{kolecki} y~\cite{kolecki2} cubren el material fundacional
sobre tensores.
El resto del material introductorio necesario se encuentra distribuido a lo
largo de la bibliografía.

\medskip

Los temas tratados en esta monografía forman parte de la bibliografía
conocida sobre álgebra lineal y álgebra abstracta.
Lo que motivó su escritura fue la percepción de que una misma idea organizadora --- la
preservación de estructura bajo transformación --- reaparece de forma
recurrente tanto en la ciencia pura como en la aplicada.
Recoger esa unidad en un solo volumen es el propósito de esta obra.
A lo largo del texto se emplea la palabra \textit{invariante} para expresar
esta idea en sentido amplio: aquello que permanece cuando cambian las
coordenadas, las representaciones o el sustrato físico.
El lector familiarizado con el álgebra advertirá que el término posee también
significados técnicos más precisos --- subespacios invariantes, invariantes de
similaridad, invariantes topológicos --- que no deben confundirse con este uso
de carácter estructural y transversal.

\medskip

En la edición, el código \LaTeX{} y el refinamiento estilístico se empleó un
modelo de lenguaje de inteligencia artificial.

\medskip

Esta monografía presenta la estructura abstracta que subyace a todo algoritmo
y sistema de cómputo autónomo, y que por tanto permanece invariante ante
cualquier cambio en el estado de la tecnología.
Existe una contraparte de orientación más física y centrada en el hardware,
organizada en torno a un computador de los años sesenta --- the Apollo Guidance Computer --- que es más específica a un momento en la historia.
Para el contexto completo, véase:

\medskip

{\small
\url{https://jportillo34.github.io/ApolloGuidanceComputer/}
}

\clearpage

\makeatletter
\@mainmattertrue
\makeatother
\pagenumbering{arabic}

\chapter{Transformaciones Lineales, Matrices y Composición}

\[
\Big\langle
\big\langle K, +_g, \cdot \big\rangle,\;
\big\langle E, +_E \big\rangle,\;
\circ
\Big\rangle
\]

La terna anterior define un \emph{espacio vectorial $E$} sobre el \emph{cuerpo $K$} \cite{cohen}. Los elementos de $K$ se llaman \emph{escalares}, y los elementos de $E$ se llaman \emph{vectores}. La operación
\(
\circ : K \times E \to E
\)
es la multiplicación por escalares, y $+_E$ es la suma de vectores. En lo que sigue, los símbolos $\alpha_1,\dots,\alpha_k$ denotan siempre escalares en $K$, y los símbolos $x_1,\dots,x_k$ denotan siempre vectores en $E$.

\medskip

Informalmente, un espacio vectorial es un conjunto de vectores que pueden sumarse entre sí y escalarse por elementos de un cuerpo.

\bigskip
\hrule
\medskip

\textit{Ejemplo (Espacio coordenado $n$-dimensional, \cite{greub}).}

Sea $\Gamma$ un cuerpo. Consideremos el conjunto
\[
\Gamma^n = \Gamma \times \dots \times \Gamma
\]
de $n$-uplas $(\xi^1,\dots,\xi^n)$, con $\xi^i \in \Gamma$. Definimos la suma por

\[
(\xi^1,\dots,\xi^n)+(\eta^1,\dots,\eta^n)=(\xi^1+\eta^1,\dots,\xi^n+\eta^n)
\]

y la multiplicación por escalar por

\[
\lambda(\xi^1,\dots,\xi^n)=(\lambda\xi^1,\dots,\lambda\xi^n),\quad \lambda \in \Gamma.
\]

Con estas operaciones, $\Gamma^n$ es un espacio vectorial sobre $\Gamma$, llamado el \emph{$n$-espacio sobre $\Gamma$}. En particular, $\Gamma$ es en sí mismo un espacio vectorial sobre $\Gamma$, donde la multiplicación por escalar coincide con la multiplicación del cuerpo.

\medskip
\hrule
\bigskip

Dados vectores $x_1,\dots,x_k \in E$ y escalares $\alpha_1,\dots,\alpha_k \in K$, una expresión de la forma
\[
\alpha_1 \circ x_1 +_E \alpha_2 \circ x_2 +_E \dots +_E \alpha_k \circ x_k
\]
se llama \emph{combinación lineal}.

\medskip

Un subconjunto $S \subset E$ se llama \emph{sistema de generadores} (o \emph{conjunto generador}) de $E$ si todo vector $x \in E$ puede escribirse como combinación lineal de elementos de $S$ \cite{greub}. Es decir, para todo $x \in E$ existen vectores $s_1,\dots,s_k \in S$ y escalares $\alpha_1,\dots,\alpha_k \in K$ tales que
\[
x = \alpha_1 \circ s_1 +_E \dots +_E \alpha_k \circ s_k.
\]

\medskip

Una familia finita de vectores $\{x_1,\dots,x_k\} \subset E$ se llama \emph{linealmente independiente} si la relación
\[
\alpha_1 \circ x_1 +_E \dots +_E \alpha_k \circ x_k = 0
\]
implica
\[
\alpha_1 = \dots = \alpha_k = 0.
\]
Si existe tal relación no trivial (es decir, con al menos un coeficiente $\alpha_i\neq 0$), la familia se llama \emph{linealmente dependiente}.

\medskip

Un conjunto finito de vectores $\{b_1,\dots,b_n\} \subset E$ se llama \emph{base} del espacio vectorial $E$ si es a la vez un sistema de generadores de $E$ y linealmente independiente.

\medskip

Equivalentemente, $\{b_1,\dots,b_n\}$ es una base de $E$ si todo vector $x \in E$ puede escribirse \emph{de manera única} como combinación lineal
\[
x = \alpha_1 \circ b_1 +_E \alpha_2 \circ b_2 +_E \dots +_E \alpha_n \circ b_n,
\quad \alpha_1,\dots,\alpha_n \in K.
\]

Una base proporciona así un conjunto mínimo y no redundante de bloques constructivos para todo el espacio. Una vez fijada una base, todo vector queda completamente determinado por sus coeficientes respecto a ella, por lo que podemos decir que la base genera el espacio.

\medskip

Si el espacio vectorial $E$ admite una base finita, el número de vectores en cualquier base de $E$ se llama la \emph{dimensión} de $E$ y se denota por $\dim E$.

\bigskip
\hrule
\medskip

% CORRECCIÓN 1: La «Observación sobre coordenadas» aparecía originalmente aquí y usaba
% el término «isomorfismo lineal» antes de definirlo formalmente. La observación se ha
% trasladado a continuación de la definición de isomorfismo, donde el término está
% disponible para el lector.

Sean $E$ y $G$ espacios vectoriales sobre el mismo cuerpo $K$. Una función
\(
\psi : E \to G
\)
se llama \emph{transformación lineal} (también conocida como \textit{aplicación lineal}) si preserva las combinaciones lineales \cite{iribarren}:
\[
\psi(\alpha_1 \circ_E x_1 + \dots + \alpha_k \circ_E x_k)
=
\alpha_1 \circ_G \psi(x_1) + \dots + \alpha_k \circ_G \psi(x_k).
\]

Las transformaciones lineales preservan la estructura algebraica de los espacios vectoriales con independencia de cualquier elección de base.

\medskip

\textit{Definición (Isomorfismo).}
Sean $E$ y $F$ espacios vectoriales sobre el mismo cuerpo $K$.
Una transformación lineal $\varphi : E \to F$ se llama \emph{isomorfismo}
si para todo $y \in F$ existe un único $x \in E$ tal que $\varphi(x) = y$,
y $\varphi(x_1) = \varphi(x_2)$ implica $x_1 = x_2$.
En este caso, decimos que $E$ y $F$ son \emph{isomorfos}.

\medskip

Un isomorfismo preserva toda la estructura lineal:
la suma, la multiplicación por escalares, la independencia lineal, la dimensión
y la validez de las relaciones lineales.
Por tanto, los espacios vectoriales isomorfos son estructuralmente idénticos.

\bigskip
\hrule
\medskip

% CORRECCIÓN 1 (continuación): La observación sobre coordenadas se traslada aquí,
% tras la definición de isomorfismo.

\textit{Observación (Sobre coordenadas e isomorfismo \cite{halmos}).}

Sea $E$ un espacio vectorial de dimensión finita sobre $K$ con $\dim E = n$.
Una vez elegida una base $B=\{b_1,\dots,b_n\}$, todo vector $x\in E$
queda representado de manera única por su columna de coordenadas
\[
[x]_B \in K^n.
\]
Esta correspondencia define un isomorfismo lineal
\[
E \longrightarrow K^n.
\]

\medskip

Así, desde el punto de vista puramente algebraico, todo espacio vectorial
$n$-dimensional sobre $K$ es estructuralmente indistinguible de $K^n$.

\medskip

Sin embargo, esta identificación depende de la base elegida.
Bases distintas producen representaciones en coordenadas distintas.
Las propiedades intrínsecas del espacio son precisamente aquellas
que permanecen invariantes bajo tales cambios de sistema de referencia.

\medskip

Por esta razón, tratamos los espacios vectoriales como entidades abstractas,
independientes de cualquier realización coordenada particular.
Las coordenadas son herramientas de cálculo; la estructura es el objeto de estudio.

\medskip
\hrule
\bigskip

Definimos el \emph{núcleo} de $\psi$ como
\[
    \ker(\psi)=\{x\in E:\psi(x)=0\},
\]
es decir, el conjunto de vectores de $E$ que la transformación envía al vector cero de $G$. Definimos la \emph{imagen} de $\psi$ como
\[
\operatorname{Im}(\psi)=\{\psi(x):x\in E\}\subseteq G,
\]

es decir, la imagen es el conjunto de valores en $G$ que $\psi$ efectivamente alcanza.

\medskip

Si $E$ es de dimensión finita, entonces
\[
\dim E = \dim \ker(\psi) + \dim \operatorname{Im}(\psi),
\]
relación conocida como el teorema de la dimensión (o teorema rango-nulidad).

\medskip

Decimos que una transformación lineal $\psi$ es \emph{inyectiva} si $\psi(x)=\psi(y)\Rightarrow x=y$ (equivalentemente, $\ker(\psi)=\{0\}$), y \emph{sobreyectiva} si para todo $g\in G$ existe $x\in E$ tal que $\psi(x)=g$ (equivalentemente, $\operatorname{Im}(\psi)=G$).

\medskip

Sean $B_E=\{e_1,\dots,e_n\}$ y $B_G=\{g_1,\dots,g_m\}$ bases de $E$ y $G$. Cada imagen $\psi(e_i)$ puede escribirse de manera única como
\(
\psi(e_i) = \alpha_{1i} \circ g_1 +_G \dots +_G \alpha_{mi} \circ g_m.
\)
Los escalares $\alpha_{ji}$ forman una tabla rectangular $m\times n$
\[
A =
\begin{bmatrix}
\alpha_{11} & \alpha_{12} & \dots & \alpha_{1n} \\
\alpha_{21} & \alpha_{22} & \dots & \alpha_{2n} \\
\vdots & \vdots & & \vdots \\
\alpha_{m1} & \alpha_{m2} & \dots & \alpha_{mn}
\end{bmatrix},
\]
y su forma abreviada
\[
A=(\alpha_{ij}),
\]
se llama la \emph{matriz de $\psi$ respecto a las bases elegidas}.

\medskip

Sea $A=(\alpha_{ij})$ una matriz $m \times n$ sobre un cuerpo $K$.

Para cada $i=1,\dots,m$, la $i$-ésima fila
\[
(\alpha_{i1},\dots,\alpha_{in})
\]
se llama \textit{vector fila} de $A$.
Se considera naturalmente como un elemento de $K^n$.

Para cada $j=1,\dots,n$, la $j$-ésima columna
\[
(\alpha_{1j},\dots,\alpha_{mj})^T
\]
se llama \textit{vector columna} de $A$.
Se considera naturalmente como un elemento de $K^m$.

\medskip

Cuando $A$ representa una transformación lineal $\psi:E\to G$
respecto a bases elegidas,
la $j$-ésima columna de $A$ es precisamente el vector de coordenadas
de $\psi(e_j)$ en la base de $G$.

\medskip

Este hecho se sigue directamente de la linealidad.
Sea $B_E=\{e_1,\dots,e_n\}$ la base elegida de $E$.
Todo vector $x\in E$ admite una representación única
\[
x=\alpha_1 e_1+\cdots+\alpha_n e_n.
\]
Como $\psi$ preserva las combinaciones lineales,
\[
\psi(x)=\alpha_1\psi(e_1)+\cdots+\alpha_n\psi(e_n).
\]

Así, la acción de $\psi$ sobre un vector arbitrario queda completamente
determinada por sus valores sobre los vectores de la base.
Si las coordenadas de cada $\psi(e_j)$ respecto a la base de $G$
se disponen como la $j$-ésima columna de una matriz $A$, entonces el vector
de coordenadas de $\psi(x)$ se obtiene formando la misma combinación lineal
de esas columnas con coeficientes $\alpha_1,\dots,\alpha_n$.
En coordenadas,
\[
[\psi(x)]_{B_G}=A[x]_{B_E}.
\]
Así, una transformación lineal queda completamente determinada por las imágenes
de los vectores de la base, y la matriz registra precisamente esas imágenes
en forma coordinada.

\medskip

Las matrices son representaciones concretas de transformaciones lineales una vez elegidas las bases.

\medskip

Las transformaciones lineales también pueden aplicarse sucesivamente.
Comprender cómo se combinan sus representaciones coordenadas
conduce de manera natural al producto de matrices.

\medskip

El producto de matrices es la expresión coordenada de la \emph{composición de aplicaciones lineales}.

\medskip

\textit{Derivación (Producto de matrices a partir de la composición).}

\medskip

Sean
\[
S : K^p \longrightarrow K^n,
\qquad
T : K^n \longrightarrow K^m
\]
transformaciones lineales con matrices (respecto a bases elegidas)
\[
B \in M_{n\times p}(K),
\qquad
A \in M_{m\times n}(K).
\]

Para un vector $x \in K^p$, el producto matriz--vector representa
la aplicación de la transformación lineal correspondiente:
\[
S(x)=Bx,
\qquad
T(y)=Ay.
\]

La composición $T\circ S : K^p \to K^m$ satisface por tanto
\[
(T\circ S)(x)=T(S(x))=A(Bx).
\]

Dado que una matriz representa de manera única una transformación lineal en coordenadas,
debe existir una matriz $C$ tal que
\[
(T\circ S)(x)=Cx
\quad \text{para todo } x\in K^p.
\]
Por definición, esta matriz es el producto $C=AB$, y queda determinada por la condición
\[
(AB)x=A(Bx).
\]
Escribiendo esta relación en coordenadas se obtiene que las entradas de $C=AB$
vienen necesariamente dadas por
\[
(AB)_{ij}
=
\sum_{k=1}^{n} A_{ik} B_{kj},
\qquad
1\le i\le m,\; 1\le j\le p.
\]

Así, la regla habitual del producto de matrices es la expresión coordenada de la composición de aplicaciones lineales.

\medskip

% CORRECCIÓN 2: La sección de cambio de base reabrió originalmente la justificación
% del producto matriz-vector después de haberla establecido ya mediante la composición.
% Se presenta ahora como una aplicación distinta —interpretar los vectores como estados
% de un sistema— en lugar de una segunda derivación de la misma regla de producto.

En muchos casos prácticos, un vector no representa meramente un elemento abstracto de un
espacio vectorial, sino el \emph{estado} de un sistema expresado respecto a una base elegida.
Llamaremos a dicha elección de base un \emph{sistema de referencia}.
Este enfoque no añade ningún contenido algebraico nuevo, pero motiva una pregunta natural:
¿qué le ocurre a la representación coordenada de un estado cuando el sistema de referencia cambia?

\medskip

% CORRECCIÓN 3: La introducción de «estado» y «sistema de referencia» incluye ahora
% una frase de enlace explícita que conecta el lenguaje abstracto usado anteriormente
% con la terminología aplicada que se introduce aquí.

Sea $E$ un espacio vectorial de dimensión finita y sea
\[
B=\{e_1,\dots,e_n\}
\]
una base de $E$. Todo vector $x\in E$ puede escribirse de manera única como
\[
x = \alpha_1 \circ e_1 +_E \dots +_E \alpha_n \circ e_n.
\]

Los escalares $\alpha_1,\dots,\alpha_n$ son las \emph{coordenadas} del estado $x$ respecto al sistema de referencia $B$. Escribir estas coordenadas como vector columna,

\[
[x]_B =
\begin{bmatrix}
\alpha_1\\
\vdots\\
\alpha_n
\end{bmatrix},
\]
es una conveniencia notacional: codifica la acción de las combinaciones lineales respecto a la base elegida.

\medskip

Supongamos ahora que se elige un segundo sistema de referencia
\[
B'=\{e'_1,\dots,e'_n\}
\]
para el mismo espacio $E$. El vector $x$ es el mismo objeto abstracto, pero sus coordenadas respecto a $B'$ forman un vector columna distinto $[x]_{B'}$.

\medskip

El paso de $[x]_B$ a $[x]_{B'}$ no es arbitrario. Está determinado por una transformación lineal
\[
T : K^n \longrightarrow K^n,
\]
cuya matriz depende únicamente de la relación entre las dos bases. Así,
\[
[x]_{B'} = A\,[x]_B,
\]
donde $A$ es la matriz del cambio de sistema de referencia.
Esto es una instancia directa del producto matriz--vector $[\psi(x)]_{B_G} = A[x]_{B_E}$
derivado anteriormente: aquí la «transformación» es precisamente el reetiquetado de coordenadas
inducido al pasar de $B$ a $B'$.

\medskip

En coordenadas, cada fila de la matriz representa la expresión coordenada de una forma lineal sobre $K^n$ (formalizaremos esta identificación al tratar los espacios duales), y el producto matriz--vector consiste precisamente en aplicar cada una de estas formas al vector de coordenadas. Concretamente, si $A$ es una matriz $m\times n$ y $x$ es un vector columna en $K^n$, entonces la $i$-ésima componente de $Ax$ es el valor de la forma lineal dada por la $i$-ésima fila de $A$ aplicada a $x$; esto está definido únicamente cuando $x$ tiene $n$ componentes, es decir, solo cuando el número de columnas de $A$ es igual al número de filas de $x$.

\medskip

La condición de que el número de columnas de $A$ sea igual al número de componentes de $[x]_B$ no es una convención; expresa el hecho de que el codominio de una aplicación lineal debe coincidir con el dominio de la siguiente.

\medskip

En este sentido, los vectores columna representan estados, las matrices representan transformaciones entre sistemas de referencia, y el producto de matrices representa cambios sucesivos de coordenadas. Las reglas computacionales habituales son, por tanto, consecuencias de la estructura abstracta introducida anteriormente.

\medskip

% CORRECCIÓN 4: Se ha eliminado la frase duplicada sobre sistemas no lineales.
% Se conserva únicamente la versión más informativa.

Las transformaciones lineales constituyen así la columna vertebral estructural de muchos modelos cuantitativos.
Incluso en sistemas donde dominan los efectos no lineales, los cálculos se organizan típicamente
en torno a aplicaciones y composiciones repetidas de aplicaciones lineales. El marco lineal
proporciona por tanto un lenguaje computacional universal, incluso cuando los fenómenos modelados
no son en sí mismos lineales.

\medskip

% CORRECCIÓN 5: Se añade un párrafo de enlace que conecta la discusión sobre
% composición/matrices con el espacio dual, motivando el nuevo concepto a partir
% de lo que el lector ya ha visto.

La derivación del producto de matrices puso de manifiesto que un vector fila actúa sobre un vector columna para producir un escalar: cada fila de una matriz define una regla que asigna un número a cualquier elemento de $K^n$, y dicha regla es lineal. Esto sugiere que los vectores fila y los vectores columna son objetos de naturaleza fundamentalmente distinta. Precisar esta distinción requiere introducir el \emph{espacio dual}.

\medskip

El conjunto de todas las transformaciones lineales de $E$ en el cuerpo base $K$ se denota por
\[
E^* = L(E,K)
\]
y se llama el \emph{espacio dual} de $E$. Un elemento $x^* \in E^*$ se llama \emph{forma lineal} (o \emph{funcional lineal}). Así,
\[
x^* : E \to K
\]
satisface
\[
x^*(\alpha_1 \circ_E x_1 +_E \dots +_E \alpha_k \circ_E x_k)
=
\alpha_1 x^*(x_1) + \dots + \alpha_k x^*(x_k).
\]

Las formas lineales asignan valores escalares a los vectores preservando la estructura lineal.

\medskip

Si $E$ es de dimensión finita con base $B_E=\{e_1,\dots,e_n\}$, entonces todo $x^*\in E^*$ queda unívocamente determinado por
\[
x^*(e_1),\dots,x^*(e_n).
\]
Estos escalares forman un vector fila
\[
[\,x^*(e_1)\;\dots\;x^*(e_n)\,].
\]

En cambio, los vectores de $E$ se representan como vectores columna. Estos dos objetos viven en espacios distintos: $E$ y $E^*$.

\medskip

Dada una matriz
\[
A = (\alpha_{ij}) \in M_{m \times n}(K),
\]
su \emph{transpuesta} es la matriz
\[
A^T = (\alpha_{ji}) \in M_{n \times m}(K),
\]
obtenida intercambiando filas y columnas.

\medskip

Así, la $i$-ésima fila de $A$ pasa a ser la $i$-ésima columna de $A^T$,
y la $j$-ésima columna de $A$ pasa a ser la $j$-ésima fila de $A^T$.

\medskip

La operación de transposición intercambia los roles de filas y columnas.
En particular, si una matriz representa una aplicación lineal
\( \psi : K^n \to K^m \),
entonces su transpuesta tiene tamaño invertido y puede interpretarse,
a nivel puramente de tablas, como un intercambio de los índices del dominio y el codominio.

Los vectores columna representan elementos de un espacio vectorial $E$, mientras que los vectores fila representan formas lineales, elementos del espacio dual $E^*$. Son objetos algebraicamente distintos, aunque puedan tener el mismo número de componentes.

\medskip

En muchas situaciones, sin embargo, un vector se utiliza para definir una forma lineal. Esto requiere estructura adicional (Capítulo II).

Un vector columna representa un elemento de $K^n$, un vector fila representa un elemento de $(K^n)^*$, y una matriz $m\times n$ representa una aplicación lineal
\[
K^n \xrightarrow{\;\psi\;} K^m.
\]

Un producto de la forma
\[
[\,x^*\,]\,A
\]
está definido porque corresponde a la composición
\[
K^n \xrightarrow{\;\psi\;} K^m \xrightarrow{\;x^*\;} K.
\]

Las dimensiones encajan precisamente porque el codominio de la primera aplicación coincide con el dominio de la segunda.

\medskip

En cambio, un producto de dos vectores fila intentaría componer
\[
K^n \to K
\quad\text{con}\quad
K^n \to K,
\]
lo que carece de sentido: la salida de la primera aplicación no pertenece al dominio de la segunda. El producto queda por tanto indefinido por necesidad.

\medskip

Las dimensiones de una matriz codifican la compatibilidad de dominios y codominios. Solo puede multiplicarse matrices cuando las aplicaciones lineales correspondientes pueden componerse.

\medskip

\textit{Ejemplo numérico concreto.}

Sea
\[
x^*(x,y)=3x+2y
\]
una forma lineal sobre $K^2$. Respecto a la base canónica, queda representada por el vector fila
\[
[\,3\;\;2\,].
\]

Sea
\[
A=
\begin{bmatrix}
2 & 3\\
5 & 4
\end{bmatrix}
\]
la matriz de una transformación lineal
\(
\psi:K^2\to K^2.
\)

El producto
\[
[\,3\;\;2\,]A
=
[\,3\;\;2\,]
\begin{bmatrix}
2 & 3\\
5 & 4
\end{bmatrix}
=
[\,19\;\;17\,]
\]
está definido porque representa la composición
\[
K^2 \xrightarrow{\;\psi\;} K^2 \xrightarrow{\;x^*\;} K.
\]

El resultado es de nuevo un vector fila, correspondiente a la forma lineal
\[
(x,y)\longmapsto 19x+17y.
\]

\medskip

En cambio, el producto
\[
[\,3\;\;2\,]\,[\,a\;\;b\,]
\]
no está definido porque intentaría componer dos aplicaciones de la forma
\[
K^2 \to K,
\]
lo cual es imposible: la salida de la primera aplicación es un escalar, no un vector de $K^2$.

\bigskip

\textit{Resumen conceptual.} Las estructuras introducidas en este capítulo ---espacios vectoriales, transformaciones lineales, bases, sistemas de coordenadas y espacios duales--- son todas facetas del mismo marco lineal subyacente. Las matrices codifican transformaciones lineales respecto a una base elegida, los vectores columna codifican estados respecto a un sistema de referencia, y los sistemas de ecuaciones lineales son simplemente expresiones coordenadas de estas transformaciones que preguntan si un vector pertenece a la imagen de una aplicación. Algoritmos como la eliminación gaussiana, si bien son esenciales para el cálculo, no crean estructuras nuevas: exploran las consecuencias a nivel coordenado de las relaciones lineales abstractas ya presentes. En conjunto, estos conceptos ilustran cómo las definiciones algebraicas abstractas dan lugar a las herramientas computacionales y geométricas habituales utilizadas en ciencias e ingeniería.

\section*{Notas y Observaciones}

\textit{Nota (¿Por qué la linealidad?, \cite{oxford}).}

\textit{Una razón elemental para el carácter fundacional del álgebra lineal dentro de las matemáticas es que, bajo condiciones adecuadas, las funciones pueden aproximarse localmente mediante funciones lineales. Este hecho, junto con la tendencia de los sistemas naturales a situarse cerca del mínimo de su energía potencial, explica la prevalencia de la linealidad en los fenómenos científicos.}

\medskip

\textit{Por ahora, diremos informalmente que una situación exhibe linealidad si podemos identificar una salida numérica y una entrada numérica tales que la primera es proporcional a la segunda.}

\bigskip
\hrule
\medskip

\textit{Nota: Sumas directas y descomposición de espacios vectoriales \cite{halmos}}

\medskip

Dados dos espacios vectoriales $U$ y $V$ sobre el mismo cuerpo $K$,
su \emph{suma directa} es el espacio vectorial
\[
U \oplus V
=
\{(x,y) \mid x \in U,\; y \in V\},
\]
con la suma y la multiplicación por escalar definidas componente a componente.

\medskip

Los subespacios
\[
\{(x,0) \mid x \in U\}
\quad\text{y}\quad
\{(0,y) \mid y \in V\}
\]
pueden identificarse con $U$ y $V$, respectivamente.
Todo elemento de $U \oplus V$ admite una descomposición única
\[
(x,y) = (x,0) + (0,y),
\]
de modo que $U \oplus V$ se obtiene combinando dos componentes independientes.

\medskip

Más en general, se dice que un espacio vectorial $W$ es la suma directa de los subespacios
$U$ y $V$ si todo $w \in W$ puede escribirse de manera única como
\[
w = u + v,
\qquad
u \in U,\; v \in V.
\]
En este caso se escribe $W = U \oplus V$.

\medskip

La noción de suma directa expresa un principio estructural:
un espacio puede descomponerse en partes independientes cuyas contribuciones
se combinan aditivamente sin interacción.

\medskip

Esta construcción debe distinguirse de las extensiones algebraicas.
Por ejemplo, los números complejos $\mathbb{C}$ forman un espacio vectorial
de dimensión dos sobre $\mathbb{R}$ y pueden identificarse, como espacios vectoriales,
con $\mathbb{R} \oplus \mathbb{R}$ mediante
\[
a + ib \;\longleftrightarrow\; (a,b).
\]
Sin embargo, esta identificación no preserva la multiplicación.
El producto de números complejos acopla las dos componentes,
introduciendo una interacción que no está presente en la suma directa.

\medskip

Así, aunque algunas extensiones pueden verse como sumas directas al nivel de
espacios vectoriales, la estructura algebraica adicional puede introducir relaciones
entre componentes que van más allá de la combinación lineal.

\medskip

La suma directa proporciona la manera fundamental de ensamblar componentes lineales independientes.
Construcciones posteriores, como el producto tensorial, introducirán un tipo distinto
de combinación en la que los componentes interactúan de forma multiplicativa en lugar de aditiva.

\medskip
\hrule
\bigskip

% CORRECCIÓN 6: Las notas están ahora ordenadas temáticamente para proporcionar un hilo
% de lectura más claro. El orden es:
%   (a) ¿Por qué la linealidad? — motivación
%   (b) Sumas directas — descomposición estructural, extiende el álgebra del capítulo
%   (c) Sistemas lineales — reinterpretación a nivel coordenado, puente hacia el cálculo
%   (d) La matriz como función — definición fundacional, ancla la visión abstracta
%   (e) Producto de matrices generalizado — abstracción algebraica de la multiplicación
%   (f) Polinomios — ejemplo de dimensión infinita, muestra la linealidad más allá de coordenadas

\bigskip
\hrule
\medskip

\textit{Nota (Sistemas de ecuaciones lineales).}

Operacionalmente, un sistema de ecuaciones lineales de la forma
\[
\begin{cases}
a_{11} x_1 + \dots + a_{1n} x_n = b_1\\
\vdots\\
a_{m1} x_1 + \dots + a_{mn} x_n = b_m.
\end{cases}
\]

se trata como una colección de ecuaciones a resolver en las incógnitas
\(
x_1,\dots,x_n.
\)

\medskip

Desde el punto de vista abstracto, este sistema no es más que la expresión coordenada de una transformación lineal.

\medskip

En efecto, sea
\[
\psi : K^n \longrightarrow K^m
\]
la transformación lineal cuya matriz respecto a las bases canónicas es
\[
A =
\begin{bmatrix}
a_{11} & \dots & a_{1n}\\
\vdots & & \vdots\\
a_{m1} & \dots & a_{mn}
\end{bmatrix}.
\]

\medskip

Escribiendo

\[
x =
\begin{bmatrix}
x_1\\
\vdots\\
x_n
\end{bmatrix}
\quad\text{y}\quad
b =
\begin{bmatrix}
b_1\\
\vdots\\
b_m
\end{bmatrix},
\]
el sistema anterior se expresa de manera equivalente como la ecuación única
\[
\psi(x) = b,
\quad\text{o, en coordenadas,}\quad
A x = b.
\]

Así, un sistema de ecuaciones lineales especifica un vector
\(
b \in K^m
\)
y pregunta si pertenece a la imagen de una transformación lineal
\(
\psi : K^n \to K^m.
\)
Las cuestiones de existencia, unicidad o multiplicidad de soluciones corresponden a propiedades geométricas de esta aplicación, como la inyectividad y la sobreyectividad, y son independientes de cualquier algoritmo particular utilizado para calcularlas.

\medskip

Algoritmos como la eliminación gaussiana son métodos para encontrar un vector
\(
x \in K^n
\)
\textit{que satisfaga}
\(
A x = b
\)
respecto a un sistema de referencia (base) elegido. Desde la perspectiva abstracta, son simplemente procedimientos para explorar la imagen y el núcleo de una transformación lineal en coordenadas. Es decir, la eliminación gaussiana aprovecha la estructura de la aplicación lineal codificada por la matriz \(A\) para determinar si existe solución, si es única y cómo expresar todas las soluciones posibles, operando enteramente en el sistema de coordenadas elegido.

\medskip
\hrule
\bigskip

\bigskip
\hrule
\medskip

\textit{Nota (La matriz como función, \cite{iribarren}).}

\medskip

Sea $S\neq\varnothing$ un conjunto y sean $m,n\ge 1$.
Una \emph{matriz $m\times n$ con entradas en $S$} es una función
\[
M:\{1,\dots,m\}\times\{1,\dots,n\}\longrightarrow S.
\]
Escribimos $a_{ij}=M(i,j)$ y abreviamos $M$ como $M=(a_{ij})$ (o $M=(a_{ij})_{m\times n}$) cuando no hay lugar a confusión.

\medskip

\textit{Caso particular.}
Cuando $S=K$ es un cuerpo, el conjunto de todas las matrices $m\times n$ sobre $K$ se denota por $M_{m\times n}(K)$.

\medskip
\hrule
\bigskip

\bigskip
\hrule
\medskip

\textit{Nota (Producto de matrices generalizado, \cite{knuth3}).}

Existe una abstracción útil del producto de matrices que hace explícito qué propiedades algebraicas se utilizan realmente. Si $A$ es una matriz $m\times n$ y $B$ es una matriz $n\times s$, y si $\circ$ y $\bullet$ son operaciones binarias tales que $\circ$ es asociativa, puede definirse el \emph{producto de matrices generalizado}
\[
A\,\genfrac{}{}{0pt}{}{\circ}{\bullet}\,B
\]
como la matriz $m\times s$ cuyas entradas son
\[
C_{ij} = (A_{i1}\bullet B_{1j}) \circ (A_{i2}\bullet B_{2j}) \circ \cdots \circ (A_{in}\bullet B_{nj}),
\qquad 1\le i\le m,\; 1\le j\le s.
\]
El producto matricial ordinario se obtiene cuando $\circ$ es $+$ y $\bullet$ es $\cdot$.

\medskip
\hrule
\bigskip

\bigskip
\hrule
\medskip

\textit{Ejemplo (Los polinomios como construcción lineal libre, \cite{knuth2}).}

Sea $K$ un cuerpo. Consideremos el conjunto
\[
K[x] = \left\{
\sum_{k=0}^{n} a_k x^k
\;\middle|\;
n \in \mathbb{N},\; a_k \in K
\right\},
\]
cuyos elementos se llaman \emph{polinomios} en el símbolo formal $x$ con coeficientes en $K$.

\medskip

El símbolo $x$ no toma ningún valor numérico; sirve únicamente como generador de expresiones formales.
La suma y la multiplicación por escalar se definen coeficiente a coeficiente:
\[
\left( \sum_{k=0}^{n} a_k x^k \right)
+_{K[x]}
\left( \sum_{k=0}^{m} b_k x^k \right)
=
\sum_{k=0}^{\max(n,m)} (a_k +_K b_k) x^k,
\]
\[
\alpha \circ
\left( \sum_{k=0}^{n} a_k x^k \right)
=
\sum_{k=0}^{n} (\alpha \cdot a_k) x^k.
\]

\medskip

Con estas operaciones, $K[x]$ es un espacio vectorial sobre $K$.

\medskip

La familia
\[
\{ 1, x, x^2, x^3, \dots \}
\]
es linealmente independiente y genera $K[x]$. Por tanto, forma una base de $K[x]$.

\medskip

Todo polinomio puede escribirse por tanto de manera única como combinación lineal
\[
u(x)
=
a_0 \circ 1
+_{K[x]}
a_1 \circ x
+_{K[x]}
\dots
+_{K[x]}
a_n \circ x^n,
\qquad a_k \in K.
\]

\medskip

El espacio vectorial $K[x]$ es de dimensión infinita, pues su base contiene infinitos elementos.

\medskip

Conceptualmente, $K[x]$ es el espacio vectorial libremente generado por las potencias formales
\[
\{ x^k \}_{k \ge 0}.
\]
No existen relaciones entre estos generadores más allá de las exigidas por los axiomas de espacio vectorial.

\medskip

La multiplicación de polinomios introduce una operación adicional sobre $K[x]$,
\[
x^i \cdot x^j = x^{i+j},
\]
que es compatible con la estructura lineal pero no es necesaria para la definición del espacio vectorial en sí.

\medskip

Así, los polinomios proporcionan un ejemplo canónico de estructura lineal generada por símbolos formales. El marco lineal aparece con independencia de la geometría, las coordenadas o cualquier interpretación numérica.
La multiplicación enriquece esta estructura, pero la linealidad es primaria.

\medskip
\hrule
\bigskip
\chapter{Geometría a partir de la Estructura Lineal}

En el Capítulo I, los espacios vectoriales se introdujeron como objetos puramente algebraicos.
Un vector era un elemento abstracto de un conjunto dotado de suma y multiplicación por escalares,
y las transformaciones lineales eran funciones que preservaban dicha estructura.
No se asumió ni se requirió ninguna noción de longitud, ángulo, distancia u ortogonalidad.

\medskip

El álgebra por sí sola nos permite decidir si los vectores son independientes,
si un conjunto genera un espacio,
o si una transformación lineal es inyectiva o sobreyectiva.
Sin embargo, no permite formular preguntas geométricas.
En muchas situaciones, no basta con saber \emph{qué} vectores son posibles;
también necesitamos saber cuán \emph{grandes} son, cuán \emph{próximos} están entre sí,
o cuán \emph{alineados} se encuentran.
Estas nociones se expresan mediante longitud, distancia y ángulo.
Sin estructura adicional, el lenguaje de los espacios vectoriales no distingue
entre «mayor» y «menor», ni permite hablar de perpendicularidad o de la
«mejor aproximación» de un vector mediante otros
(como en las proyecciones y los mínimos cuadrados).
No existe ninguna manera intrínseca, dentro de un espacio vectorial desnudo,
de comparar el tamaño de dos vectores o de determinar si apuntan en direcciones
perpendiculares.

\medskip

La geometría comienza únicamente cuando se impone estructura adicional.

\medskip

Conviene recordar ahora cómo se representa habitualmente el espacio geométrico
en coordenadas.
Dados conjuntos $A_1,\dots,A_n$, su \emph{producto cartesiano}
(también llamado \emph{producto de conjuntos}) es
\[
A_1\times\cdots\times A_n = \{(x_1,\dots,x_n)\mid x_i\in A_i\}.
\]
En particular, el espacio coordenado $\mathbb{R}^n$ es el producto cartesiano de
$n$ copias de $\mathbb{R}$.
Esta construcción por sí sola no es todavía geometría: proporciona un
\emph{modelo} conveniente en el que los «puntos» se representan mediante
$n$-uplas de números reales.
La geometría emerge únicamente después de especificar una regla adicional que
nos permita medir y comparar dichas $n$-uplas.

\medskip

Un tema central de este documento es que la geometría no reemplaza a la
estructura lineal, sino que la enriquece.
Los espacios vectoriales introducidos anteriormente permanecen inalterados;
lo que cambia es que se equipan ahora con una regla que permite a los vectores
interactuar de un modo que produce cantidades escalares.

\medskip

Informalmente, esta regla toma dos vectores y devuelve un escalar.
A partir de esta única operación emergen de manera natural nociones como
longitud, ángulo y distancia.

\medskip

Para precisar esta interacción, introducimos la noción de aplicación bilineal.

\medskip

Sean $E$ y $G$ espacios vectoriales sobre el mismo cuerpo $K$.
Una aplicación
\[
B : E \times G \longrightarrow K
\]
se llama \emph{bilineal} si es lineal en cada argumento por separado.
Es decir, para $y \in G$ fijo, la aplicación
\[
x \longmapsto B(x,y)
\]
es una forma lineal sobre $E$, y para $x \in E$ fijo, la aplicación
\[
y \longmapsto B(x,y)
\]
es una forma lineal sobre $G$.

\medskip

Explícitamente, para todo $x_1,x_2 \in E$, $y_1,y_2 \in G$, y todo escalar
$\alpha \in K$,
\[
B(x_1 + x_2, y) = B(x_1, y) + B(x_2, y), \qquad
B(\alpha x, y) = \alpha B(x, y),
\]
\[
B(x, y_1 + y_2) = B(x, y_1) + B(x, y_2), \qquad
B(x, \alpha y) = \alpha B(x, y).
\]

\medskip

Cuando $E = G$, dicha aplicación se llama \emph{forma bilineal} sobre $E$.

\medskip

A este nivel no se asume ninguna interpretación geométrica.
Una forma bilineal es simplemente una regla que asigna un escalar a un par
ordenado de vectores, sujeta a la linealidad en cada argumento.
No tiene por qué ser simétrica, positiva, ni estar relacionada con ninguna
noción de longitud o ángulo.

\medskip

No obstante, las formas bilineales desempeñan un papel estructural fundamental.
Proporcionan una manera sistemática de asociar vectores con formas lineales
y de construir cantidades escalares a partir de pares de vectores.

\medskip

Las nociones geométricas familiares emergen únicamente después de imponer
condiciones adicionales sobre la forma bilineal.
Estos casos especiales se introducirán en breve.

% CORRECCIÓN 3: La restricción a R se introduce ahora aquí, inmediatamente después
% de las formas bilineales y antes de cualquier mención a positividad o productos
% internos. Anteriormente aparecía mucho más tarde, después de que el producto
% interno ya se había usado implícitamente.

\medskip

Hasta este punto, el cuerpo escalar $K$ ha sido arbitrario.
La estructura lineal pura no requiere la capacidad de comparar escalares.
La geometría, sin embargo, depende de la positividad: expresiones como
$\langle x,x\rangle > 0$ y desigualdades de la forma $r>0$
presuponen una relación de orden.
Por esta razón, restringimos ahora la atención al sistema de los números reales
\[
\Big\langle \mathbb{R}, +, \cdot, \le \Big\rangle,
\]
visto como cuerpo ordenado.
Es precisamente el orden en $\mathbb{R}$ el que permite que las longitudes,
las distancias y la ortogonalidad adquieran significado geométrico.

\medskip

Especializamos ahora las formas bilineales introducidas anteriormente.

\medskip

Sea $E$ un espacio vectorial sobre $\mathbb{R}$.
Una \emph{forma bilineal} sobre $E$ es una función
\[
B : E \times E \longrightarrow \mathbb{R}
\]
que es lineal en cada argumento por separado.
Es decir, para todo $x,x',y,y' \in E$ y todo $\lambda \in \mathbb{R}$,
\[
B(x+x',y)=B(x,y)+B(x',y), \qquad
B(\lambda x,y)=\lambda B(x,y),
\]
y análogamente en el segundo argumento.

\medskip

Las formas bilineales surgen de manera natural a partir de las transformaciones lineales.
Si
\[
\psi : E \longrightarrow E^*
\]
es una aplicación lineal, entonces la fórmula
\[
B(x,y) = \psi(x)(y)
\]
define una forma bilineal sobre $E$.
Recíprocamente, toda forma bilineal determina tal aplicación.
Por tanto, las formas bilineales codifican una manera controlada en que los
vectores pueden interactuar para producir escalares.

\medskip

A este nivel no está presente ninguna geometría.
Una forma bilineal puede ser degenerada, asimétrica o indefinida.
Permite que los vectores interactúen, pero todavía no mide nada.

\medskip

La geometría aparece cuando la forma bilineal satisface restricciones adicionales.

% CORRECCIÓN 1 y 2: El producto interno se define ahora antes de cualquier
% expresión coordenada que lo involucre. La sección sobre la transpuesta, que
% anteriormente interrumpía el camino de las formas bilineales al producto interno,
% ha sido reubicada después de que el producto interno esté establecido. La
% expresión coordenada del producto interno sigue a la definición, no la precede.

\medskip

Un \emph{producto interno} sobre un espacio vectorial $E$ sobre $\mathbb{R}$
es una forma bilineal
\[
\langle \cdot,\cdot \rangle : E \times E \longrightarrow \mathbb{R}
\]
que es simétrica, positiva y no degenerada.
Estas propiedades garantizan que
\[
\langle x,x \rangle > 0 \quad \text{para todo } x \neq 0.
\]

\medskip

Una vez fijada dicha regla, los vectores adquieren propiedades medibles.
Dos vectores pueden declararse ortogonales.
A un vector se le puede asignar una longitud.
La distancia entre dos puntos puede definirse en términos de los vectores que
los unen.
Ninguno de estos conceptos existe antes de la introducción de esta estructura adicional.

\medskip

Como en el Capítulo I, las coordenadas no desempeñan ningún papel fundamental.
Las cantidades geométricas deben ser independientes de cualquier sistema de
referencia particular.
Una fórmula que cambia al cambiar la base describe una representación,
no un objeto geométrico.

\medskip

Solo después de elegir una base aparecen las expresiones coordenadas familiares.
Las sumas de cuadrados, los productos escalares y las ecuaciones cuadráticas
surgen como manifestaciones a nivel coordenado de la estructura subyacente,
no como definiciones de ella.

\medskip

\textit{Expresión coordenada del producto interno.}

Sea $(E,\langle\cdot,\cdot\rangle)$ un espacio con producto interno de dimensión
finita y sea $\{e_1,\dots,e_n\}$ una base ortonormal.
Todo vector $x\in E$ puede escribirse de manera única como
\[
x=\sum_{i=1}^n x_i e_i,
\qquad
y=\sum_{i=1}^n y_i e_i.
\]

Usando la bilinealidad y la condición de ortonormalidad
$\langle e_i,e_j\rangle=\delta_{ij}$, obtenemos
\[
\langle x,y\rangle
=
\sum_{i=1}^n x_i y_i.
\]

Así, en una base ortonormal el producto interno abstracto queda representado
por el familiar \emph{producto escalar} de vectores coordenados en $\mathbb{R}^n$.

\medskip

Esta observación explica el papel del producto de matrices.
Si una transformación lineal $\psi:E\to G$ está representada por una matriz
$A=(a_{ij})$ y un vector $x$ tiene vector coordenado $x=(x_1,\dots,x_n)^T$,
entonces la $i$-ésima componente del producto $Ax$ es
\[
(Ax)_i=\sum_{j=1}^n a_{ij}x_j.
\]

Cada componente es por tanto el producto escalar de la $i$-ésima fila de $A$
con el vector $x$.
El producto de matrices puede así entenderse como la evaluación sistemática de
productos internos entre filas de una matriz y columnas de otra.

\medskip

En este sentido, la regla algebraica familiar para multiplicar matrices no es
una convención computacional arbitraria: admite una interpretación geométrica
natural como expresión coordenada de interacciones medidas por un producto interno.

\medskip

Este punto de vista explica por qué objetos como circunferencias, esferas y
elipsoides pueden describirse mediante ecuaciones algebraicas y, sin embargo,
conservan un significado geométrico invariante bajo cambios de coordenadas.
Las ecuaciones dependen de la base elegida; la geometría no.

\medskip

El objetivo es por tanto mostrar cómo la geometría emerge de la estructura
lineal una vez que se especifica una interacción adecuada entre los vectores.

\medskip

% CORRECCIÓN 2 (continuación): La sección sobre la transpuesta y la transformación
% dual se ubica ahora aquí, después de definir el producto interno. Esto preserva
% la derivación algebraica de la transpuesta como intrínseca (sin necesidad de
% producto interno), a la vez que permite que la observación sobre la dualidad
% interna aparezca de manera natural en el contexto de un producto interno ya establecido.

Aunque la geometría ha comenzado a emerger a través de la interacción bilineal,
el mecanismo algebraico subyacente sigue siendo la transformación lineal.
Para comprender cómo interactúa la estructura geométrica con las aplicaciones
lineales, examinamos una construcción asociada de manera natural a toda
transformación lineal.

\medskip

El objeto central de este trabajo es la transformación lineal. Antes de
introducir nociones geométricas como longitud, ángulo y distancia, hacemos
explícita una construcción fundamental asociada a toda transformación lineal:
su transpuesta.

\medskip

Sean $E$ y $G$ espacios vectoriales sobre el mismo cuerpo $K$, y sea
\[
\psi : E \longrightarrow G
\]
una transformación lineal. Recordemos que los espacios duales $E^*$ y $G^*$
consisten en todas las formas lineales sobre $E$ y $G$, respectivamente.

\medskip

Para cada $g^* \in G^*$, la composición
\[
g^* \circ \psi : E \longrightarrow K
\]
es una forma lineal sobre $E$. Esto define una aplicación
\[
\psi^* : G^* \longrightarrow E^*,
\qquad
\psi^*(g^*) = g^* \circ \psi.
\]

\medskip

La aplicación $\psi^*$ es lineal y se llama la \emph{transpuesta}
(o transformación dual) de $\psi$. La dirección de la flecha se invierte:
\[
E \xrightarrow{\;\psi\;} G
\quad\Longrightarrow\quad
G^* \xrightarrow{\;\psi^*\;} E^*.
\]

\medskip

Esta construcción es intrínseca: depende únicamente de la transformación
lineal $\psi$ y no requiere ninguna elección de base, coordenadas o producto interno.

\medskip

Cuando se eligen bases para $E$ y $G$, y se eligen las bases duales
correspondientes para $E^*$ y $G^*$, la transformación lineal $\psi$ queda
representada por una matriz $A$, y la transformación dual $\psi^*$ queda
representada por la matriz transpuesta $A^T$.
Así, la transposición matricial familiar es la expresión coordenada de la
transformación dual intrínseca $\psi^*$.

\medskip

El producto interno hace más que asignar longitudes.
Establece una identificación canónica entre el espacio vectorial y su dual.

\medskip

Para cada vector $x \in E$, definimos la aplicación
\[
\Phi(x) : E \longrightarrow \mathbb{R},
\qquad
\Phi(x)(y) = \langle x,y \rangle.
\]

Por bilinealidad, $\Phi(x)$ es una forma lineal sobre $E$.
Obtenemos así una aplicación
\[
\Phi : E \longrightarrow E^*.
\]

\medskip

Si el producto interno es no degenerado, esta aplicación es inyectiva.
En dimensión finita, como $\dim E = \dim E^*$, es por tanto un isomorfismo.

\medskip

Así, una vez fijado un producto interno, cada vector determina una forma lineal
única, y toda forma lineal proviene de un vector único.
La distinción entre $E$ y $E^*$ desaparece a nivel estructural.

\medskip

Es precisamente esta identificación la que permite considerar la transpuesta
de una transformación lineal como un operador sobre el propio $E$.
Cuando se introduce la geometría, la dualidad se vuelve interna.

\medskip

Una vez fijado un producto interno, los vectores adquieren longitud.
La norma de un vector $x \in E$ se define por
\[
\|x\| = \sqrt{\langle x,x \rangle}.
\]

\medskip

Así, el espacio vectorial $E$ introducido en el Capítulo I queda ahora dotado
de estructura geométrica adicional:
\[
\Big\langle
\langle \mathbb{R},+,\cdot\rangle,\;
\langle E,+\rangle,\;
\circ,\;
\langle\cdot,\cdot\rangle
\Big\rangle.
\]

El producto interno no es necesario para la linealidad en sí, pero permite
introducir nociones geométricas como longitud, ángulo, ortogonalidad y
transformaciones adjuntas.

\medskip

\textit{Expansión algebraica de la norma.}

Para cualesquiera vectores $x,y \in E$, el cuadrado de la norma de su suma satisface
\[
\|x+y\|^2
=
\langle x+y, x+y \rangle.
\]
Por bilinealidad y simetría del producto interno,
\[
\langle x+y, x+y \rangle
=
\langle x,x \rangle
+
\langle x,y \rangle
+
\langle y,x \rangle
+
\langle y,y \rangle
=
\|x\|^2
+
2\langle x,y \rangle
+
\|y\|^2.
\]

Por tanto,
\[
\|x+y\|^2
=
\|x\|^2
+
\|y\|^2
+
2\langle x,y \rangle.
\]

Esta identidad se cumple en todo espacio con producto interno.
Expresa cómo el término de interacción $\langle x,y \rangle$ mide la
desviación respecto a la aditividad pura de los cuadrados de las longitudes.

\medskip

El producto interno satisface también la desigualdad
\[
|\langle x,y \rangle|
\le
\|x\|\,\|y\|,
\]
conocida como la \emph{desigualdad de Cauchy--Schwarz}.
Como consecuencia,
\[
\|x+y\|^2
\le
\|x\|^2
+
\|y\|^2
+
2\|x\|\,\|y\|
=
(\|x\|+\|y\|)^2.
\]

Tomando raíces cuadradas se obtiene la \emph{desigualdad triangular}
\[
\|x+y\|
\le
\|x\|
+
\|y\|.
\]

\medskip

La \textit{identidad de Pitágoras} surge como el caso especial
$\langle x,y \rangle = 0$,
en el que el término de interacción se anula y
\[
\|x+y\|^2
=
\|x\|^2
+
\|y\|^2.
\]

\medskip

La ortogonalidad también adquiere significado.
Se dice que dos vectores $x,y \in E$ son ortogonales si
\[
\langle x,y \rangle = 0.
\]

\medskip

El orden de los vectores es irrelevante.
La simetría del producto interno implica
\[
\langle x,y \rangle = \langle y,x \rangle,
\]
de modo que la ortogonalidad es una relación mutua.

\medskip

El vector cero es ortogonal a todo vector de $E$.
En efecto, para cualquier $v \in E$,
\[
\langle 0,v \rangle
=
\langle 0\cdot v, v \rangle
=
0\langle v,v \rangle
=
0.
\]
Más aún, el vector cero es el único con esta propiedad.

\medskip

Así, la ortogonalidad es precisamente la condición bajo la cual el término
de interacción se anula.
En ese caso,
\[
\|x+y\|^2 = \|x\|^2 + \|y\|^2.
\]
Los cuadrados de las longitudes se vuelven aditivos exactamente cuando los
vectores no interactúan.

\medskip

% CORRECCIÓN 4: El interludio sobre vectores propios se ha trasladado a aquí,
% después de establecer la norma, la ortogonalidad y el adjunto. Anteriormente
% aparecía entre la identidad de Pitágoras y la sección sobre distancia,
% interrumpiendo la cadena norma -> distancia -> topología en su punto más crítico.

\textit{Interludio: Direcciones invariantes}

Hasta este punto, las transformaciones lineales se han tratado algebraicamente:
una aplicación $\psi : E \to E$ envía vectores a vectores preservando la
estructura lineal.
Todavía no se ha atribuido ningún significado geométrico a esta acción.

\medskip

Una vez introducido un producto interno, sin embargo, la acción de una
transformación lineal puede examinarse geométricamente.
Dado un vector $x \in E$, la imagen $\psi(x)$ puede diferir de $x$ de dos
maneras independientes: su longitud puede cambiar y su dirección puede cambiar.

\medskip

Para un vector genérico, ambos efectos ocurren simultáneamente.
La mayoría de los vectores son estirados y rotados.

\medskip

Se produce un fenómeno notable cuando existe un vector no nulo $x \in E$
tal que
\[
\psi(x) = \lambda x
\]
para algún escalar $\lambda$.
En este caso, la transformación actúa sobre $x$ mediante un reescalado puro.
La dirección de $x$ se preserva.

\medskip

Tal vector se llama \emph{vector propio} de $\psi$, y el escalar $\lambda$
es su \emph{valor propio} correspondiente.
Los vectores propios identifican direcciones en el espacio que son
geométricamente invariantes bajo la transformación.

\medskip

Esta invariancia es una propiedad intrínseca de la propia transformación lineal.
Cambiar la base puede alterar la representación de $\psi$, pero no puede
destruir ni crear vectores propios.

\medskip

Bajo la identificación $E \cong E^*$ inducida por el producto interno, la
transpuesta $\psi^*$ corresponde a un operador lineal sobre $E$, llamado
el \emph{adjunto} de $\psi$.

\medskip

Cuando una transformación lineal proviene de una forma bilineal simétrica (o,
equivalentemente, de un operador autoadjunto), sus vectores propios
correspondientes a valores propios distintos son ortogonales.
Estas direcciones definen un sistema de coordenadas preferido dictado por la
geometría, no por una elección arbitraria.

\medskip

Es en este sentido que los vectores propios sirven como ejes principales de
los elipsoides asociados a las formas cuadráticas.
La diagonalización es el acto de alinear las coordenadas con las direcciones
invariantes ya presentes en la estructura.

\medskip

% CORRECCIÓN 4 (continuación): Con los vectores propios establecidos aquí,
% la sección sobre distancia y topología sigue sin interrupción.

\textit{Distancia y emergencia de la geometría.}

Una vez disponible un producto interno sobre un espacio vectorial $E$, este
determina una norma, y a partir de la norma obtenemos una noción de distancia
entre vectores,
\[
d(x,y)=\|x-y\|.
\]

En este punto no se introduce ninguna estructura nueva: la distancia es
simplemente una reformulación de datos algebraicos ya presentes.

\medskip

Esta distancia nos permite describir la proximidad. Para un vector $x \in E$
y un número real $r > 0$, consideramos el conjunto
\[
B(x,r) = \{\, y \in E \mid d(x,y) < r \,\}.
\]
Dicho conjunto se llama \textit{bola abierta} centrada en $x$ con radio $r$.
Geométricamente, consiste en todos los vectores que se encuentran más cerca
de $x$ que el umbral prescrito.

\medskip

La desigualdad estricta es esencial. Garantiza que todo punto interior a la
bola admite puntos cercanos adicionales, de modo que no se incluye ninguna
frontera.
Este «margen de holgura» es lo que en última instancia nos permite hablar de
continuidad y deformación.

\medskip

La colección de todas las bolas abiertas determina qué subconjuntos de $E$
deben llamarse abiertos: un conjunto es abierto si, alrededor de cada uno de
sus puntos, contiene alguna bola abierta.
De este modo, la distancia genera una noción de apertura sin axiomas adicionales.

% CORRECCIÓN 5: Una frase de transición prepara ahora al lector para la
% definición axiomática de espacio topológico, conectándola con la intuición
% geométrica construida justo antes, en lugar de introducir los axiomas
% abruptamente.

\medskip

Esta noción de apertura puede capturarse mediante un pequeño conjunto de
axiomas que preservan sus propiedades esenciales y permiten extender el concepto
mucho más allá del contexto métrico.
La estructura resultante es lo que se conoce como una \textit{topología} sobre $E$.

\medskip

Un \textit{espacio topológico} es un par ordenado $(X, \tau)$, donde $X$ es un
conjunto y $\tau$ es una colección de subconjuntos de $X$.
Para que $\tau$ sea una \textit{topología}, debe satisfacer los siguientes axiomas:

\begin{enumerate}
    \item $\emptyset \in \tau$ y $X \in \tau$.
    \item La unión de cualquier colección de conjuntos en $\tau$ también pertenece a $\tau$.
    \item La intersección de cualquier número \textit{finito} de conjuntos en $\tau$
    también pertenece a $\tau$.
\end{enumerate}

Los elementos de $\tau$ se llaman \textit{conjuntos abiertos}.

\medskip

Así, la topología no aparece aquí como una teoría independiente, sino como una
consecuencia estructural de la métrica inducida por el producto interno.
La geometría da lugar a la distancia, la distancia da lugar a los entornos,
y los entornos dan lugar a la estructura topológica.

\medskip

\textit{Interpretación geométrica.}

En $\mathbb{R}^n$ dotado de su producto interno estándar, las bolas abiertas
corresponden a las esferas familiares de la geometría euclídea.
Normas más generales deforman estas bolas en formas elipsoidales, reflejando
un escalado anisotrópico a lo largo de distintas direcciones del espacio.
En todos los casos, cada punto
\[
x = (x_1,\dots,x_n)
\]
es un vector del espacio vectorial ambiente, y los objetos geométricos
considerados son subconjuntos de ese espacio definidos por relaciones de
distancia entre sus vectores.

\medskip

Una vez que los objetos quedan representados como vectores de este modo, las
estructuras geométricas introducidas anteriormente ---normas, distancias, bolas
y elipsoides--- se convierten en herramientas para organizarlos y compararlos.

\medskip

Sus posiciones relativas, medidas mediante una métrica elegida, codifican
similitud.
Esta métrica no tiene por qué ser euclídea: en la práctica puede diseñarse
para reflejar invariancias propias del dominio o incluso aprenderse a partir
de datos.
Los entornos de significado se describen entonces mediante bolas abiertas,
y el aprendizaje puede interpretarse como el descubrimiento de una organización
geométrica y topológica subyacente a los datos.

\medskip

Una vez introducida una norma, ciertos conjuntos geométricos emergen como
consecuencias directas de la estructura y no requieren hipótesis adicionales.

\medskip

Sea $(E,\langle\cdot,\cdot\rangle)$ un espacio con producto interno y sea
$\|x\|=\sqrt{\langle x,x\rangle}$ la norma asociada.
Para un vector fijo $x_0\in E$ y un escalar $r>0$, la \emph{bola cerrada}
de radio $r$ centrada en $x_0$ se define como
\[
B_r(x_0)=\{\,x\in E \mid \|x-x_0\|\le r\,\}.
\]

La \emph{esfera} correspondiente es la frontera de este conjunto,
\[
S_r(x_0)=\{\,x\in E \mid \|x-x_0\|= r\,\}.
\]

\textit{Observación (Descripción coordenada en $\mathbb{R}^n$):}
Cuando $E=\mathbb{R}^n$ está dotado de su producto interno estándar y se
utilizan las coordenadas estándar $x=(x_1,\dots,x_n)$, las mismas definiciones
intrínsecas toman la familiar forma de «suma de cuadrados».
Para una bola centrada en el origen,
\[
B_r(0)=\left\{x\in\mathbb{R}^n\;\middle|\;x_1^2+\cdots+x_n^2\le r^2\right\},
\]
y su esfera frontera es
\[
S_r(0)=\left\{x\in\mathbb{R}^n\;\middle|\;x_1^2+\cdots+x_n^2=r^2\right\}.
\]
Esto es simplemente la expresión coordenada de
$\|x\|=\sqrt{x_1^2+\cdots+x_n^2}$.

\medskip

Estas definiciones son intrínsecas y no dependen de ninguna elección de
coordenadas o base.

\medskip

\textit{Observación (Simetría):}
En el caso euclídeo $E=\mathbb{R}^n$ con su producto interno estándar, la
bola es invariante bajo transformaciones ortogonales.
Si $Q\in O(n)$, entonces $\|Qx\|=\|x\|$, y por tanto
\[
Q\bigl(B_r(0)\bigr)=B_r(0).
\]
Más en general, toda isometría de un espacio con producto interno transforma
bolas en bolas.

\medskip

\textit{Observación (Secciones en $\mathbb{R}^n$):}
En $\mathbb{R}^n$, fijar la última coordenada $x_n=t$ muestra que la
sección transversal de una bola de radio $r$ por el hiperplano $x_n=t$ es
una bola $(n-1)$-dimensional de radio $\sqrt{r^2-t^2}$.
Esta observación contiene el contenido geométrico que subyace al método
habitual de «integrar por secciones» para calcular volúmenes $n$-dimensionales.

\medskip

\textit{Observación (Volumen de bolas euclídeas):}
En $\mathbb{R}^n$, el volumen $n$-dimensional de una bola euclídea escala como
una constante multiplicada por $r^n$:
\[
\operatorname{Vol}_n\bigl(B_r(0)\bigr)=\omega_n\,r^n,
\]
donde $\omega_n=\operatorname{Vol}_n\bigl(B_1(0)\bigr)$ es el volumen de la
bola unidad. Una fórmula cerrada es
\[
\omega_n=\frac{\pi^{n/2}}{\Gamma\!\left(\frac{n}{2}+1\right)}.
\]
En particular, $\omega_1=2$, $\omega_2=\pi$ y $\omega_3=\frac{4\pi}{3}$.

\medskip

\textit{Observación (Un efecto en alta dimensión):}
Aunque $B_1(0)\subset \mathbb{R}^n$ contiene siempre muchos puntos, su
volumen $n$-dimensional satisface $\omega_n\to 0$ cuando $n\to\infty$.
Esto ilustra un fenómeno básico de la geometría en alta dimensión: conjuntos
que parecen grandes en baja dimensión pueden ocupar una fracción despreciable
del volumen ambiente al aumentar la dimensión.

\medskip

Más en general, sea $B:E\times E\to \mathbb{R}$ una forma bilineal simétrica
definida positiva. El conjunto
\[
\mathcal{E}=\{\,x\in E \mid B(x-x_0,x-x_0)\le r^2\,\}
\]
se llama \emph{elipsoide}.
Así, un elipsoide es la bola unidad asociada a una forma bilineal, o
equivalentemente, a un producto interno modificado sobre $E$.

\medskip

Cuando la forma bilineal $B$ es simétrica y definida positiva, existe un
único operador lineal autoadjunto $A:E\to E$ tal que
\[
B(x,y)=\langle Ax,y\rangle
\]
para todo $x,y\in E$.

\medskip

El operador $A$ codifica la geometría del elipsoide.
Como $A$ es autoadjunto sobre un espacio con producto interno de dimensión
finita, admite una base ortonormal de vectores propios.
Si $\{e_1,\dots,e_n\}$ es tal base y $A e_i=\lambda_i e_i$ con
$\lambda_i>0$, entonces la desigualdad definidora
\[
B(x-x_0,x-x_0)\le r^2
\]
se convierte en
\[
\sum_{i=1}^n \lambda_i\, \xi_i^2 \le r^2,
\]
donde $\xi_i=\langle x-x_0,e_i\rangle$ son las coordenadas en esta
base propia.

% CORRECCIÓN 6: La observación de que la diagonalización «revela las direcciones
% invariantes ya presentes» aparecía dos veces en sucesión cercana. Se conserva
% la primera aparición (aquí, en la derivación del elipsoide) ya que concluye el
% argumento algebraico. La reiteración redundante que aparecía unas líneas más
% tarde, justo antes de la fórmula coordenada, ha sido eliminada.

\medskip

Así, los ejes principales del elipsoide son precisamente los vectores propios
de $A$, y los valores propios determinan el escalado a lo largo de esas
direcciones.
La diagonalización no crea los ejes; simplemente revela las direcciones
invariantes ya presentes en la estructura.

\medskip

En consecuencia, cuando se elige una base en la que la forma bilineal es
diagonal, esta definición invariante se reduce a una expresión coordenada
familiar.
En $\mathbb{R}^n$, se obtiene
\[
\mathcal{E}
=
\left\{
(x_1,\dots,x_n)\in\mathbb{R}^n
\;\middle|\;
\sum_{i=1}^n \frac{(x_i-a_i)^2}{h_i^2} \le r^2
\right\}.
\]

Escrita por componentes, esta desigualdad es
\[
\left(\frac{x_1-a_1}{h_1}\right)^2
+
\left(\frac{x_2-a_2}{h_2}\right)^2
+
\cdots
+
\left(\frac{x_n-a_n}{h_n}\right)^2
\le r^2.
\]

La ecuación depende de las coordenadas elegidas; el elipsoide en sí, no.

\medskip

\emph{Observación (Casos degenerados).}
La naturaleza geométrica del conjunto definido anteriormente depende de los
coeficientes $h_1,\dots,h_n$.
Si $h_1=\cdots=h_n$, el elipsoide se reduce a una esfera $n$-dimensional.
Si uno o más de los parámetros $h_i$ se anulan, la forma cuadrática
definidora se vuelve degenerada y el elipsoide colapsa en un objeto de
dimensión inferior: un disco, un segmento o, en el caso extremo, un punto.

\medskip

Desde el punto de vista estructural, estas degeneraciones corresponden a una
pérdida de rango en la forma cuadrática o bilineal asociada.
Así, los cambios en la estructura algebraica se traducen directamente en
cambios en la dimensión geométrica.

\medskip

\emph{Observación.}
A lo largo de esta discusión, el espacio ambiente $E=\mathbb{R}^n$ es un
espacio vectorial en el sentido del Capítulo I.
Cada elemento
\[
x=(x_1,\dots,x_n)
\]
es por tanto un vector de $E$.

\medskip

Sin embargo, los conjuntos definidos por estas expresiones son \emph{subconjuntos}
de $E$ seleccionados mediante restricciones cuadráticas.
No son espacios vectoriales en sí mismos, pues no son cerrados bajo la suma
ni bajo la multiplicación por escalares.

\bigskip

\textit{Resumen conceptual.}
Los espacios vectoriales proporcionan el escenario.
Las transformaciones lineales describen los cambios de estado admisibles.
La geometría aparece únicamente cuando se permite a los vectores interactuar
mediante una regla bilineal estructurada que produce escalares.
A partir de esta interacción emergen longitudes, ángulos, distancias y formas
geométricas.
La transición del álgebra a la geometría no es por tanto un salto, sino un
enriquecimiento controlado de la estructura.
El libro \cite{iribarren2} es una referencia excelente sobre la topología de
los espacios métricos.

\section*{Notas y Observaciones}

\textit{Perspectivas.}

En aplicaciones contemporáneas como el aprendizaje de representaciones,
los elementos de un conjunto (palabras, imágenes o señales) se proyectan en
espacios vectoriales de alta dimensión, donde las relaciones geométricas
codifican similitud e información estructural.
Estas construcciones proporcionan realizaciones concretas de ideas abstractas
del álgebra lineal.

\medskip

Los sistemas computacionales modernos ofrecen otro ejemplo ilustrativo.
Las plataformas dedicadas a la \emph{observabilidad} ---es decir, a la medición
sistemática del estado interno de un sistema en ejecución--- recopilan grandes
familias de cantidades numéricas como latencia, rendimiento, consumo de
memoria, tasas de error y utilización de recursos.

\medskip

En un instante de tiempo fijo, estas mediciones pueden ensamblarse en un vector
\[
x(t) = (x_1(t),\dots,x_n(t)) \in \mathbb{R}^n,
\]
cuyas coordenadas representan características observables del sistema.
A medida que evoluciona el tiempo, el sistema determina una trayectoria en este
espacio lineal ambiente, de modo que la monitorización y el análisis pueden
interpretarse geométricamente como el estudio de trayectorias en un espacio
vectorial de dimensión finita.

\medskip

Las consideraciones operacionales imponen típicamente cotas sobre cada observable,
expresadas mediante desigualdades de la forma
\[
a_i \le x_i \le b_i,
\]
que reflejan limitaciones físicas o rangos de operación aceptables.
Geométricamente, estas restricciones definen un producto cartesiano de
intervalos, es decir, un hiperrectángulo $n$-dimensional (o, tras
normalización, un hipercubo).
El comportamiento normal del sistema puede verse por tanto como confinado a
una región acotada del espacio ambiente, mientras que el comportamiento
anómalo corresponde a trayectorias que se alejan de dicha región.

\medskip

Conviene subrayar que esta construcción es un modelo matemático idealizado:
en los sistemas prácticos, la colección de observables puede variar con el
tiempo, por lo que la dimensión $n$ debe considerarse fija únicamente de manera
local o dentro de una representación elegida.
No obstante, el marco de los espacios lineales proporciona un lenguaje natural
para describir estados, restricciones y evolución.

\medskip

De este modo, los sistemas de observabilidad ilustran un tema recurrente del
álgebra lineal: un espacio ambiente plano sirve como escenario universal sobre
el cual los fenómenos significativos se caracterizan mediante subconjuntos
definidos geométricamente.
Las restricciones de norma dan lugar a esferas, los escalados anisotrópicos
producen elipsoides, y las cotas coordenadas definen hiperrectángulos.
A pesar de sus diferentes apariencias geométricas, todos ellos emergen como
subconjuntos estructurados embebidos en el mismo espacio lineal ambiente.

\bigskip
\hrule
\medskip

% CORRECCIÓN 7: La observación histórica sobre Lambert se ha trasladado a
% inmediatamente después del Interludio sobre Direcciones Invariantes en la
% sección de Notas, donde su conexión con los vectores propios y la preservación
% de ángulos es contextualmente inmediata. Anteriormente aparecía más adelante
% en las notas, separada de la discusión que la motivaba.

\textit{Observación histórica (Ángulos e invariancia).}

Antes de que los vectores propios se comprendieran como direcciones invariantes,
su detección solía ser indirecta.
En dos dimensiones, determinar si una transformación preserva una dirección se
reduce a detectar si preserva un ángulo.

\medskip

Este problema condujo de manera natural al cálculo de arcotangentes.
En el siglo XVIII, J.H.~Lambert introdujo la fracción continua
\[
\arctan(z)
=
\cfrac{z}{1 + \cfrac{z^2}{3 + \cfrac{4z^2}{5 + \cfrac{9z^2}{7 + \cdots}}}},
\]
en el marco de sus trabajos sobre mecánica orbital y fenómenos de rotación.

\medskip

Desde el punto de vista moderno, este esfuerzo analítico refleja una pregunta
geométrica: ¿rota una transformación lineal una dirección, o simplemente la
escala?
Los vectores propios proporcionan la respuesta estructural a esta pregunta,
haciendo del cálculo de ángulos algo auxiliar en lugar de fundamental.

\medskip
\hrule
\bigskip

\textit{Observación (Formas cuadráticas y el estadístico chi-cuadrado).}
Las expresiones de la forma que definen elipsoides aparecen de manera natural
en contextos que no son, a primera vista, geométricos.
Un ejemplo notable surge en la teoría de las pruebas de números aleatorios,
donde se encuentra la cantidad
\[
V
=
\frac{(Y_1 - np_1)^2}{np_1}
+
\frac{(Y_2 - np_2)^2}{np_2}
+
\cdots
+
\frac{(Y_n - np_n)^2}{np_n},
\]
conocida en estadística como el \emph{estadístico chi-cuadrado} asociado a
las cantidades observadas $Y_1,\dots,Y_n$ y los valores de referencia
$np_1,\dots,np_n$.

\medskip

Desde el presente punto de vista, esta expresión es simplemente una forma
cuadrática sobre $\mathbb{R}^n$.
Mide el cuadrado de la longitud del vector de desviación
\[
(Y_1 - np_1,\dots,Y_n - np_n)
\]
respecto a una forma bilineal simétrica definida positiva cuya matriz, en la
base estándar, es diagonal con entradas
$(np_1)^{-1},\dots,(np_n)^{-1}$.

\medskip

En consecuencia, los conjuntos definidos por desigualdades de la forma
\[
V \le c
\]
son elipsoides centrados en el punto $(np_1,\dots,np_n)$.
Los denominadores $np_i$ determinan el escalado anisotrópico a lo largo de
las direcciones coordenadas, exactamente como en el elipsoide general
\[
\sum_{i=1}^n \frac{(x_i-a_i)^2}{h_i^2} \le r^2.
\]

\medskip

Así, antes de imponer cualquier interpretación probabilística, la expresión
chi-cuadrado es ya un objeto geométrico: es la norma al cuadrado inducida
por un producto interno particular sobre $\mathbb{R}^n$.
La teoría estadística hace uso de esta geometría, pero no la crea.

\medskip

\textit{Observación (El espacio vectorial ambiente y las regiones restringidas).}
El propio espacio vectorial $E$ es el \emph{espacio ambiente} en el que viven
todos los vectores.
En coordenadas, cuando se elige una base, este espacio aparece como
$\mathbb{R}^n$, un entorno lineal plano en el que los vectores pueden sumarse
y escalarse libremente.

\medskip

Los objetos geométricos como bolas, esferas y elipsoides no reemplazan a este
espacio ambiente; son subconjuntos del mismo definidos por restricciones sobre
la norma o sobre una forma cuadrática.
Por ejemplo, la esfera
\[
S_r(x_0)=\{x\in E \mid \|x-x_0\|=r\}
\]
selecciona los vectores cuya distancia a $x_0$ es fija, mientras que un
elipsoide selecciona los vectores cuyas coordenadas satisfacen una desigualdad
cuadrática anisotrópica.

\medskip

En muchas aplicaciones modernas (el aprendizaje profundo en inteligencia
artificial, por ejemplo), los algoritmos operan sobre vectores en un espacio
ambiente de alta dimensión imponiendo tales restricciones durante el cálculo.
Los pasos de normalización pueden restringir los vectores a esferas, los
procedimientos de regularización pueden confinar los parámetros a regiones
elipsoidales, y los métodos de optimización buscan soluciones dentro de
entornos descritos por bolas.
El espacio vectorial ambiente proporciona el universo lineal en el que existen
los vectores, mientras que estos subconjuntos geométricos describen las regiones
seleccionadas por las reglas o restricciones del problema.
\chapter{Tensores como Aplicaciones Multilineales}

La linealidad expresa que una aplicación respeta la estructura algebraica de
un espacio vectorial.
Sin embargo, muchas construcciones en matemáticas dependen de funciones que
toman varios argumentos vectoriales simultáneamente y permanecen lineales en
cada uno de ellos cuando los demás se mantienen fijos. Esto conduce de manera
natural a la noción de multilinealidad.

\medskip

Sean \(E_1, E_2, \dots, E_n\) y \(F\) espacios vectoriales sobre el mismo
cuerpo \(K\).

\begin{definition}[Aplicación multilineal]
Una aplicación
\[
T : E_1 \times E_2 \times \cdots \times E_n \to F
\]
se llama \emph{multilineal} si es lineal en cada argumento por separado.
Es decir, para cada índice \(i \in \{1,\dots,n\}\) y para todo
\[
x_i, y_i \in E_i, \quad \alpha \in K,
\]
se tiene
\[
T(x_1,\dots,x_i +_{E_i} y_i,\dots,x_n)
=
T(x_1,\dots,x_i,\dots,x_n)
+_F
T(x_1,\dots,y_i,\dots,x_n),
\]
y
\[
T(x_1,\dots,\alpha \circ x_i,\dots,x_n)
=
\alpha \circ T(x_1,\dots,x_i,\dots,x_n),
\]
con todos los demás argumentos fijos.
\end{definition}

Cuando \(n=1\), una aplicación multilineal es simplemente una aplicación lineal.
Así, las aplicaciones multilineales no son un tipo de objeto nuevo, sino una
generalización directa de las transformaciones lineales.

\medskip

Las aplicaciones multilineales aparecen de manera natural en muchos contextos:
productos bilineales, emparejamientos entre vectores y formas lineales,
e interacciones de orden superior que no pueden reducirse a una única entrada.
Su estudio sistemático conduce al concepto de tensor.

\medskip

Informalmente, un tensor es un objeto diseñado para codificar comportamiento
multilineal de manera intrínseca y sin coordenadas.

\medskip

A lo largo de este documento se utilizarán dos puntos de vista equivalentes.
El primero trata los tensores directamente como aplicaciones multilineales.

\medskip

Sea \(E\) un espacio vectorial sobre \(K\), y sea \(E^*\) su espacio dual,
es decir, el espacio vectorial de las aplicaciones lineales de \(E\) en \(K\).

\begin{definition}[Tensor como aplicación multilineal]
Un \emph{tensor de tipo \((k,\ell)\)} sobre \(E\) es una aplicación multilineal
\[
T : \underbrace{E^* \times \cdots \times E^*}_{k\ \text{veces}}
\times
\underbrace{E \times \cdots \times E}_{\ell\ \text{veces}}
\;\longrightarrow\; K.
\]
\end{definition}

Tal tensor toma \(k\) formas lineales y \(\ell\) vectores como argumentos,
y produce un escalar.
Se requiere linealidad en cada argumento por separado.

\medskip

Esta definición precisa varios casos familiares:
\begin{itemize}
\item Un escalar en \(K\) puede verse como un tensor de tipo \((0,0)\).
\item Un vector en \(E\) corresponde a un tensor de tipo \((0,1)\).
\item Una forma lineal en \(E^*\) es un tensor de tipo \((1,0)\).
\item Una forma bilineal sobre \(E\) es un tensor de tipo \((0,2)\).
\item Si $E$ es de dimensión finita, una aplicación lineal $T : E \to E$ puede
      identificarse con un tensor de tipo $(1,1)$.
\end{itemize}

Así, los vectores y las matrices reaparecen como los casos de orden más bajo
de un único marco unificado.

\medskip

A este nivel de abstracción no intervienen coordenadas, tablas ni
interpretaciones geométricas.
Un tensor queda definido enteramente por cómo actúa sobre vectores y formas
lineales, y por la multilinealidad de dicha acción.

% CORRECCIÓN 1: La motivación del producto tensorial está ahora claramente
% separada de su definición. La versión anterior mezclaba motivación y solución
% en el mismo bloque, con la frase final anticipando la definición que seguía
% a continuación (circular). La motivación es ahora autocontenida: identifica la
% limitación de las aplicaciones multilineales (no existe composición natural) y
% plantea la pregunta que el producto tensorial responderá, sin nombrar la
% solución prematuramente.

\medskip

La definición anterior es conceptualmente clara, pero tiene una limitación
práctica.
Las aplicaciones lineales se componen: si $S : E \to F$ y $T : F \to G$ son
lineales, entonces $T \circ S : E \to G$ es de nuevo lineal, y su matriz es
el producto de las matrices de $T$ y $S$.
Las aplicaciones multilineales no admiten una operación análoga.
Sus salidas y entradas no encajan de manera natural como para permitir una
composición canónica, y no existe ninguna regla general para combinar dos
aplicaciones multilineales en una tercera preservando la multilinealidad.

\medskip

Las aplicaciones multilineales forman espacios vectoriales bajo las operaciones
punto a punto, de modo que pueden sumarse y escalarse.
Lo que falta es una manera de \emph{reducir} la multilinealidad a la linealidad:
sustituir un problema multilineal por uno lineal sobre un espacio
convenientemente construido.
Esto es precisamente lo que logra el producto tensorial.

\medskip

Sean \(E_1,\dots,E_n\) espacios vectoriales sobre el mismo cuerpo \(K\).
Consideremos el conjunto de todas las aplicaciones multilineales
\[
T : E_1 \times \cdots \times E_n \longrightarrow F
\]
con valores en un espacio vectorial arbitrario \(F\).
El producto tensorial se define de modo que cada tal aplicación multilineal
corresponda de manera única a una aplicación lineal definida sobre un único
espacio vectorial.

\medskip

\begin{definition}[Producto tensorial]
El \emph{producto tensorial} de los espacios vectoriales \(E_1,\dots,E_n\) es
un espacio vectorial, denotado
\[
E_1 \otimes \cdots \otimes E_n,
\]
junto con una aplicación multilineal
\[
\otimes : E_1 \times \cdots \times E_n \longrightarrow
E_1 \otimes \cdots \otimes E_n,
\]
que satisface la siguiente propiedad universal:

Para todo espacio vectorial \(F\) y toda aplicación multilineal
\[
T : E_1 \times \cdots \times E_n \longrightarrow F,
\]
existe una única aplicación lineal
\[
\widetilde{T} : E_1 \otimes \cdots \otimes E_n \longrightarrow F
\]
tal que
\[
T(x_1,\dots,x_n) = \widetilde{T}(x_1 \otimes \cdots \otimes x_n)
\]
para todo \(x_i \in E_i\).
\end{definition}

\medskip

Esta propiedad caracteriza completamente el producto tensorial.
Los elementos \(x_1 \otimes \cdots \otimes x_n\) se llaman \emph{tensores simples}.
Todo elemento del espacio producto tensorial es una combinación lineal finita
de tensores simples.

\medskip

Como consecuencia, el estudio de las aplicaciones multilineales es equivalente
al estudio de las aplicaciones lineales definidas sobre productos tensoriales.
Concretamente,
\[
\text{Aplicaciones multilineales } E_1 \times \cdots \times E_n \to F
\quad\longleftrightarrow\quad
\text{Aplicaciones lineales } E_1 \otimes \cdots \otimes E_n \to F.
\]

\medskip

Esta equivalencia es el fundamento conceptual de la teoría de tensores.
Explica por qué los tensores pueden manipularse algebraicamente, componerse
con aplicaciones lineales y representarse concretamente una vez elegidas las bases.

\medskip

Volviendo al caso de un único espacio vectorial \(E\), un tensor de tipo
\((k,\ell)\) puede verse equivalentemente como una aplicación lineal
\[
T :
\underbrace{E^* \otimes \cdots \otimes E^*}_{k}
\otimes
\underbrace{E \otimes \cdots \otimes E}_{\ell}
\;\longrightarrow\; K.
\]

\medskip

Así los tensores son aplicaciones lineales definidas sobre productos tensoriales,
y su aparente complejidad surge únicamente al introducir coordenadas.

\medskip

A diferencia de la suma directa (Capítulo~I), que combina componentes
independientes de manera aditiva, el producto tensorial codifica interacciones
entre componentes mediante la multilinealidad.

\bigskip
\hrule
\medskip

\textit{Nota:} Una vez elegidas las bases, los productos tensoriales se
convierten en espacios vectoriales de dimensión finita, y los tensores admiten
representaciones coordenadas como tablas multidimensionales. Estas
representaciones dependen de las bases elegidas y no definen el tensor en sí.

\medskip
\hrule
\bigskip

Sea $E$ un espacio vectorial de dimensión finita sobre $K$, y sea
\[
B = \{ e_1, \dots, e_n \}
\]
una base de $E$. Su base dual
\[
B^* = \{ e^1, \dots, e^n \}
\subset E^*
\]
se define por las relaciones
\[
e^i(e_j) = \delta^i_j,
\]
donde $\delta^i_j$ es la delta de Kronecker.

\medskip

Todo vector $x \in E$ admite una expansión única
\[
x = x^i \circ e_i,
\]
y toda forma lineal $\varphi \in E^*$ admite una expansión única
\[
\varphi = \varphi_i \circ e^i.
\]

Los escalares $x^i$ y $\varphi_i$ son las \emph{componentes} de $x$ y de
$\varphi$ respecto a las bases elegidas.

\medskip

Sea ahora $T$ un tensor de tipo $(k,\ell)$ sobre $E$, es decir,
\[
T : \underbrace{E^* \times \cdots \times E^*}_{k}
\times
\underbrace{E \times \cdots \times E}_{\ell}
\longrightarrow K.
\]

La acción de $T$ queda completamente determinada por sus valores sobre los
elementos de la base:
\[
T(e^{i_1}, \dots, e^{i_k}, e_{j_1}, \dots, e_{j_\ell})
= T^{i_1 \dots i_k}{}_{j_1 \dots j_\ell}.
\]

Estos escalares se llaman las \emph{componentes de $T$ respecto a la base $B$}.
Codifican el tensor por completo una vez fijada una base.

\medskip

Para argumentos arbitrarios
\[
(\varphi^1, \dots, \varphi^k, x_1, \dots, x_\ell),
\]
la multilinealidad implica la expansión
\[
T(\varphi^1, \dots, \varphi^k, x_1, \dots, x_\ell)
=
T^{i_1 \dots i_k}{}_{j_1 \dots j_\ell}
\,
\varphi^1_{i_1} \cdots \varphi^k_{i_k}
\,
x_1^{j_1} \cdots x_\ell^{j_\ell}.
\]

Así, la notación de índices es la expresión coordenada de la multilinealidad
una vez elegida una base.

\medskip

Las aplicaciones bilineales proporcionan el ejemplo computacional más simple
de multilinealidad.
Una aplicación bilineal
\[
B : V \times W \to U
\]
es lineal en cada argumento por separado. Una vez elegidas bases en los espacios
vectoriales involucrados, la aplicación queda representada por coeficientes
$t_{ij}^{\;k}$ tales que las componentes del vector de salida satisfacen
\begin{equation}
z^{k} = \sum_{i=1}^{m}\sum_{j=1}^{n} t_{ij}^{\;k}\, x^{i} y^{j}.
\end{equation}

La colección de coeficientes $(t_{ij}^{\;k})$ forma un tensor de rango tres.
Esta representación muestra que los procedimientos computacionales familiares
son expresiones coordenadas de aplicaciones multilineales. Por ejemplo, la
multiplicación de números complejos
\begin{equation}
(x_1 + i x_2)(y_1 + i y_2)
= (x_1 y_1 - x_2 y_2)
+ i(x_1 y_2 + x_2 y_1)
\end{equation}
puede interpretarse como la evaluación de dos formas bilineales determinadas
por un tensor de coeficientes $2 \times 2 \times 2$.

Desde este punto de vista, los algoritmos que manipulan tablas numéricas son
implementaciones de transformaciones tensoriales escritas en un sistema de
referencia elegido. El cálculo hace por tanto explícita la estructura
multilineal invariante subyacente a las operaciones algebraicas.

\medskip

En particular, una aplicación bilineal
\[
B : E \times F \to G
\]
entre espacios vectoriales de dimensión finita queda representada, una vez
elegidas las bases, por componentes
\[
B^k{}_{ij},
\]
de modo que para vectores $x \in E$ e $y \in F$,
\[
(B(x,y))^k = B^k{}_{ij}\, x^i y^j.
\]
Así, las tablas multidimensionales de coeficientes que aparecen en los
algoritmos numéricos son representaciones coordenadas de tensores que codifican
operaciones multilineales.

\medskip

Una vez fijada una base, las componentes
\[
T^{i_1 \dots i_k}{}_{j_1 \dots j_\ell}
\]
pueden disponerse en una tabla multidimensional de escalares.

\medskip

Cambiar la base de $E$ cambia los valores numéricos de las componentes según
reglas de transformación precisas, mientras que el tensor en sí permanece
inalterado.

\bigskip
\hrule
\medskip

\textit{Nota: Tensores mediante leyes de transformación}

La definición de tensor como aplicación multilineal es intrínseca y libre de
coordenadas. Una caracterización equivalente pone de relieve cómo se
transforman las componentes de un tensor bajo un cambio de base.

\medskip

Este punto de vista hace explícito el contenido invariante que sobrevive a los
cambios de coordenadas, que es precisamente lo que permite que las leyes
tensoriales aparezcan en física y en aplicaciones de ingeniería con
independencia de la representación elegida.

\medskip

Sean $B = \{e_i\}$ y $B' = \{e'_i\}$ dos bases de $E$, relacionadas por
\[
e'_i = A^j{}_i \, e_j,
\]
donde $A^j{}_i$ es una matriz de cambio de base invertible.
Las bases duales correspondientes satisfacen
\[
e'^i = (A^{-1})^i{}_j \, e^j.
\]

\medskip

Sea $T$ un tensor de tipo $(k,\ell)$ sobre $E$, con componentes respecto a
$B$ dadas por
\[
T^{i_1 \dots i_k}{}_{j_1 \dots j_\ell}.
\]

La matriz $A^j{}_i$ expresa los nuevos vectores de la base en términos de los
anteriores, tal como se trató en el Capítulo~I.
En consecuencia, la regla de transformación siguiente describe cómo cambia la
representación coordenada de $T$ al pasar de la base $B$ a $B'$.

\medskip

Las componentes respecto a la nueva base $B'$ vienen dadas por
\[
T'^{i_1 \dots i_k}{}_{j_1 \dots j_\ell}
=
A^{i_1}{}_{p_1}
\cdots
A^{i_k}{}_{p_k}
\,
(A^{-1})^{q_1}{}_{j_1}
\cdots
(A^{-1})^{q_\ell}{}_{j_\ell}
\,
T^{p_1 \dots p_k}{}_{q_1 \dots q_\ell}.
\]

\medskip

Los índices contravariantes (índices superiores) se transforman con la matriz
$A$, mientras que los índices covariantes (índices inferiores) se transforman
con su inversa.
Esta regla de transformación no es una hipótesis adicional.
Es una consecuencia directa de la multilinealidad y de cómo se transforman
los vectores y los vectores duales bajo un cambio de base.

\medskip

Recíprocamente, supongamos que para cada elección de base se asigna una tabla
multidimensional de escalares
\[
T^{i_1 \dots i_k}{}_{j_1 \dots j_\ell}
\]
de tal manera que estas tablas están relacionadas por la ley de transformación
anterior siempre que se cambie la base.
Entonces estas colecciones coordenadas definen una única aplicación multilineal
\[
T :
\underbrace{E^* \times \cdots \times E^*}_{k}
\times
\underbrace{E \times \cdots \times E}_{\ell}
\longrightarrow K,
\]
independiente de la base elegida.

\medskip

Así los dos puntos de vista siguientes son equivalentes:
\begin{itemize}
\item Un tensor es una aplicación multilineal.
\item Un tensor es una colección de componentes que se transforman según la
      ley de transformación tensorial.
\end{itemize}

La primera definición enfatiza la estructura intrínseca.
La segunda enfatiza la invariancia bajo cambio de sistema de referencia.
Su equivalencia explica por qué la teoría de tensores desempeña un papel
central en geometría y física: las ecuaciones tensoriales son válidas en todo
sistema de coordenadas, porque ambos lados se transforman según la misma ley.

\medskip
\hrule
\bigskip

Así:
\begin{itemize}
\item los tensores son objetos multilineales abstractos;
\item las tablas son representaciones dependientes de la base;
\item los índices registran cómo la multilinealidad se descompone en coordenadas.
\end{itemize}

Esta distinción es esencial para interpretar las implementaciones computacionales.

\medskip

Hasta este punto, los tensores se han considerado como objetos multilineales
estáticos.
Introducimos ahora la operación que les permite interactuar algebraicamente.

% CORRECCIÓN 3: La frase rota («operación. introducimos ahora...») y la
% repetición de la observación sobre la no composición han sido eliminadas.
% La transición es ahora una única frase limpia que abre la sección de contracción.

\medskip

Para combinar tensores y producir nuevos objetos multilineales, se introduce
la operación fundamental de contracción tensorial, que generaliza el producto
de matrices y explica la aparición de sumas sobre índices repetidos.

\medskip

Denotamos por
\[
\mathcal{T}^{(k,\ell)}(E)
\]
el espacio vectorial de todos los tensores de tipo $(k,\ell)$ sobre $E$.

\medskip

Sea $E$ un espacio vectorial de dimensión finita sobre $K$.
Consideremos un tensor
\[
T \in \mathcal{T}^{(k,\ell)}(E)
\]
y un tensor
\[
S \in \mathcal{T}^{(m,n)}(E).
\]

Supongamos que se selecciona un índice covariante de $T$ y un índice
contravariante de $S$.
La contracción tensorial es la operación que empareja estos dos índices mediante
la aplicación de evaluación natural
\[
E^* \times E \to K.
\]

\begin{definition}[Contracción tensorial]
Sean
\[
T \in \mathcal{T}^{(k,\ell)}(E),
\qquad
S \in \mathcal{T}^{(m,n)}(E).
\]
Eligiendo un índice covariante de $T$ y un índice contravariante de $S$, su
contracción se define componiendo estos argumentos con el emparejamiento de
evaluación natural
\[
\langle\cdot,\cdot\rangle : E^* \times E \to K.
\]
En coordenadas, esta operación aparece como una suma sobre una base.

\medskip

El tensor resultante tiene tipo
\[
(k+m-1,\;\ell+n-1).
\]
\end{definition}

\medskip

La contracción reduce el número total de índices del tensor en dos: un índice
contravariante y uno covariante se emparejan y eliminan.
No se requieren coordenadas para definir esta operación.

\medskip

Para expresar la contracción en coordenadas, debe elegirse una base.
El emparejamiento abstracto aparece entonces como una suma sobre índices repetidos.

\medskip

Fijemos una base $\{e_i\}$ de $E$ y su base dual $\{e^i\}$.

Sean
\[
T^{i}{}_{j}
\quad\text{y}\quad
S^{j}{}_{k}
\]
las componentes de dos tensores de tipo $(1,1)$.

\medskip

La contracción sobre el índice $j$ produce un nuevo tensor con componentes
\[
R^{i}{}_{k}
=
T^{i}{}_{j} \, S^{j}{}_{k}.
\]

El índice repetido $j$ indica una suma:
\[
R^{i}{}_{k}
=
\sum_{j=1}^{n}
T^{i}{}_{j} \, S^{j}{}_{k}.
\]
Esta suma no es una regla adicional.
Es la expresión coordenada del emparejamiento abstracto entre $E^*$ y $E$.

\medskip

Sean $A : E \to E$ y $B : E \to E$ aplicaciones lineales.
Cada una puede identificarse con un tensor de tipo $(1,1)$.
Respecto a una base, sus componentes se escriben como
\[
A^{i}{}_{j},
\qquad
B^{j}{}_{k}.
\]

La composición $A \circ B : E \to E$ tiene componentes
\[
(A \circ B)^{i}{}_{k}
=
A^{i}{}_{j} \, B^{j}{}_{k}.
\]

Así, el producto de matrices es exactamente la contracción tensorial entre dos
tensores de tipo $(1,1)$, y la regla de suma familiar para el producto de
matrices es un caso especial del emparejamiento abstracto entre $E^*$ y $E$.

\medskip

A nivel abstracto tenemos:
\[
\text{Contracción tensorial}
\]
que, tras elegir bases, se convierte en
\[
\text{Suma sobre índices}
\quad\longrightarrow\quad
C^{i}{}_{k} = \sum_j A^{i}{}_{j} B^{j}{}_{k}.
\]

% CORRECCIÓN 4: El párrafo que comenzaba con «Es fundamental reconocer...»
% ha sido eliminado. Su tono era marcadamente más retórico que el texto
% circundante, introducía el término no estándar «demolición de índices»,
% y su contenido (la contracción como motor del cálculo) ya se transmite
% de manera más precisa en la nota sobre Arrays Sistólicos más adelante.

% CORRECCIÓN 2: El ejemplo de la cuadrícula discreta (tabla de fuerzas) se ha
% trasladado aquí, tras definir la contracción. El lector dispone ahora de las
% herramientas algebraicas para comprender qué significan los índices de la tabla
% y cómo podrían contraerse. Anteriormente aparecía antes de la contracción,
% dejando la tabla sin ningún contexto operacional.

\medskip

En contextos computacionales y de modelado, los tensores aparecen con frecuencia
a través de tablas multidimensionales cuyos índices etiquetan parámetros
independientes del problema.
El siguiente ejemplo ilustra cómo surgen en la práctica tales tablas mientras
siguen representando datos tensoriales.

\medskip

Consideremos un sistema físico modelado sobre una cuadrícula bidimensional
discreta.
En cada punto espacial $(i,j)$ y en cada instante de tiempo $t$, se mide un
vector de fuerza
\[
F(t,i,j) \in \mathbb{R}^2
\]
con componentes $(F_x, F_y)$ respecto a un sistema de referencia espacial fijo.

\medskip

Sean:
\begin{itemize}
\item $t \in \{1,\dots,T\}$ el índice de tiempo,
\item $(i,j) \in \{1,\dots,N\} \times \{1,\dots,M\}$ el índice de posición espacial,
\item $\alpha \in \{1,2\}$ el índice de componentes de la fuerza.
\end{itemize}

Los datos recopilados pueden escribirse como una tabla
\[
F_{t i j \alpha},
\]
de forma $(T,N,M,2)$.

\medskip

Cada índice corresponde a una elección de modelado independiente: tiempo,
espacio y componentes vectoriales.
La tabla almacena valores numéricos respecto a sistemas de coordenadas elegidos.
Solo el índice de componente $\alpha$ corresponde a un índice vectorial genuino
en el sentido tensorial; los índices $(t,i,j)$ etiquetan parámetros
independientes del modelo.
Contraer el índice $\alpha$ con una forma lineal produce, en cada punto del
espacio-tiempo, una medición escalar de la fuerza en una dirección elegida.

\medskip

Abstractamente, estos datos pueden interpretarse como la representación
coordenada de una función con valores tensoriales sobre un dominio discreto.

\bigskip
\hrule
\medskip

\textit{Nota: El álgebra lineal realizada en hardware --- Arrays sistólicos
\cite{kung}}

\medskip

La interpretación del producto de matrices como contracción tensorial no es
solo una reformulación conceptual.
También determina cómo se diseña el hardware computacional moderno.

\medskip

En los grandes cálculos numéricos, la operación dominante es la evaluación
repetida de productos internos,
\[
c_i = \sum_{j=1}^{n} a_{ij} b_j,
\]
que es la expresión coordenada de una aplicación lineal aplicada a un vector.
Equivalentemente, es una contracción entre un tensor de tipo $(1,1)$ y un
tensor de tipo $(0,1)$.

\medskip

A finales de la década de 1970, H.~T.~Kung y C.~E.~Leiserson propusieron
organizar el cálculo en torno a esta operación mediante arquitecturas llamadas
\emph{arrays sistólicos}.
En lugar de transferir datos de ida y vuelta entre la memoria y un procesador
central, los datos fluyen rítmicamente a través de una cuadrícula de elementos
de procesamiento simples.
Cada elemento realiza repetidamente una única operación de la forma
\[
C \leftarrow A B + C,
\]
un paso de multiplicación-acumulación que implementa una contribución a un
producto interno.

\medskip

Matemáticamente, cada elemento de procesamiento evalúa una contracción tensorial
parcial.
A medida que los datos se propagan por el array, las contracciones sucesivas
se acumulan, y el resultado global emerge de la acción coordinada de muchas
operaciones lineales locales.

\medskip

Esta organización refleja un principio estructural: la eficiencia del hardware
se logra cuando el cálculo refleja la estructura algebraica.
Dado que las transformaciones lineales se descomponen en sumas de productos,
una arquitectura optimizada para las operaciones de multiplicación-acumulación
implementa de manera natural el álgebra lineal a gran escala.

\medskip

Las arquitecturas computacionales modernas de alto rendimiento adoptan con
frecuencia variantes de este paradigma, disponiendo muchos elementos de
procesamiento simples para evaluar operaciones de multiplicación-acumulación
en paralelo.
Lo que aparece a nivel de software como contracción tensorial o producto de
matrices se ejecuta físicamente como flujos de productos internos coordinados.

\medskip

Desde el presente punto de vista, este desarrollo no es accidental.
La contracción tensorial es la interacción fundamental entre objetos
multilineales, y el cálculo a gran escala reduce los algoritmos complejos a
evaluaciones repetidas de esta interacción.

\medskip

En este sentido, las tablas numéricas almacenadas en memoria, las sumas sobre
índices escritas en fórmulas y los datos que fluyen por circuitos de silicio
son tres realizaciones del mismo objeto matemático: la evaluación de
aplicaciones multilineales.

\medskip
\hrule
\bigskip

Los principios abstractos desarrollados anteriormente aparecen de manera
concreta en la física clásica, donde las cantidades tensoriales describen
relaciones invariantes entre observables físicas.

\medskip

\textit{Un ejemplo clásico: el tensor de inercia \cite{feynman}}

En mecánica clásica, la distribución de masa en un cuerpo rígido se codifica
mediante un tensor de segundo orden llamado \emph{tensor de inercia}.
Respecto a una base ortonormal elegida de $\mathbb{R}^3$, sus componentes
vienen dadas por
\[
I_{ij} = \int_{\Omega}
\big( \delta_{ij} \|x\|^2 - x_i x_j \big)\, dm,
\]
donde $\Omega$ es el cuerpo y $x$ es el vector de posición.

\medskip

El tensor de inercia define una relación lineal entre la velocidad angular
$\omega$ y el momento angular $L$:
\[
L_i = I_{ij} \, \omega^j.
\]

\medskip

Bajo un cambio de base determinado por una matriz invertible $A$, sus
componentes se transforman según
\[
I'_{ij}
=
A^{p}{}_{i} \,
A^{q}{}_{j} \,
I_{pq},
\]
que es precisamente la ley de transformación de un tensor de tipo $(0,2)$.

\medskip

Así, aunque las entradas numéricas de $I_{ij}$ dependen del sistema de
referencia elegido, el objeto geométrico que representa no depende de él.
El contenido físico ---la resistencia a la aceleración rotacional--- es
invariante bajo cambio de coordenadas.

\medskip

La teoría de tensores extiende por tanto el principio central introducido en
el Capítulo~I: la estructura matemática se identifica no por su representación
numérica, sino por la invariancia de sus relaciones definidoras bajo las
transformaciones admisibles.

% CORRECCIÓN 5 y 6: Las notas se reordenan para agrupar los ejemplos aplicados
% y separarlos de la perspectiva especulativa.
% Orden: Perspectivas (hardware futuro) -> observación sobre chi-cuadrado
% (puente con el Cap. II) -> cubos OLAP (aplicación en análisis de datos)
% -> Arrays sistólicos (hardware).
% El párrafo introductorio redundante que precedía a la nota sobre Arrays
% Sistólicos en el texto principal ha sido eliminado; la apertura de la propia
% nota resulta ahora suficiente.

\section*{Notas y Observaciones}

\bigskip
\hrule
\medskip

\textit{Perspectiva: Computación más allá de la Electrónica \cite{Xu}}

\medskip

Los avances recientes en arquitecturas de cómputo exploran sustratos físicos
más allá de los circuitos electrónicos tradicionales, entre ellos los sistemas
ópticos y fotónicos en los que la información se procesa mediante
interferencia, difracción y propagación de la luz.

\medskip

A pesar de sus diferencias físicas, estos sistemas están concebidos para
implementar las mismas operaciones algebraicas fundamentales que aparecen a lo
largo de este capítulo.
En particular, la tarea computacional dominante sigue siendo la evaluación de
expresiones de la forma
\[
C^{i}{}_{k} = A^{i}{}_{j} B^{j}{}_{k},
\]
es decir, la contracción tensorial.

\medskip

En estas arquitecturas emergentes, la interacción de las señales físicas
realiza de forma natural operaciones de multiplicación y acumulación en
paralelo, de modo que la evaluación de aplicaciones multilineales es ejecutada
directamente por el proceso físico subyacente, sin necesidad de instrucciones
electrónicas secuenciales.

\medskip

Sin embargo, surge una restricción no trivial debida a la codificación física
de la información.
Los sistemas ópticos no coherentes representan valores como intensidades
luminosas, que son intrínsecamente no negativas; operan, por tanto, dentro del
dominio restringido~$\mathbb{R}^+$ en lugar de sobre la recta real completa.
Dado que las contracciones tensoriales exigen pesos y activaciones que recorren
todo~$\mathbb{R}$, tales sistemas resultan algebraicamente incompletos para el
cómputo en general.
Las arquitecturas coherentes resuelven esta limitación codificando los valores
como \emph{amplitudes} de campo óptico --- magnitudes que admiten signo --- y
recuperando la salida con signo mediante interferencia en la detección.
El grupo completo $(\mathbb{R},+)$ queda así disponible, y las contracciones
tensoriales de signo arbitrario son realizables directamente en el hardware.

\medskip

Desde esta perspectiva, los avances en hardware no alteran la estructura
matemática de la computación.
Proporcionan nuevas realizaciones físicas de las mismas operaciones
invariantes: transformaciones lineales y contracciones tensoriales, ahora
extendidas a su dominio natural.
La completitud algebraica que aquí se requiere --- el paso
de~$\mathbb{R}^+$ a~$\mathbb{R}$ --- es precisamente la distinción estructural
entre un grupo multiplicativo y uno aditivo, cuya correspondencia es el objeto
del Capítulo~IV.

\medskip

Así, con independencia de si el cómputo se lleva a cabo por medios
electrónicos, ópticos o de cualquier otra naturaleza física, el marco
algebraico desarrollado aquí permanece inalterado.

\medskip
\hrule
\bigskip

\bigskip
\hrule
\medskip

\textit{Observación: Puente con el Capítulo~II --- el estadístico chi-cuadrado
como tensor de tipo $(0,2)$}

\medskip

El estadístico chi-cuadrado encontrado en el Capítulo~II proporciona un ejemplo
anterior de tensor de tipo $(0,2)$.
Recordemos que la expresión
\[
V
=
\sum_{i=1}^{n}
\frac{(Y_i - np_i)^2}{np_i}
\]
fue identificada como la norma al cuadrado inducida por una forma bilineal
simétrica definida positiva $B$ sobre $\mathbb{R}^n$, cuya matriz en la base
estándar es diagonal con entradas $(np_i)^{-1}$.
Desde el presente punto de vista, $B$ es precisamente un tensor de tipo $(0,2)$:
una aplicación multilineal
\[
B : \mathbb{R}^n \times \mathbb{R}^n \longrightarrow \mathbb{R}
\]
que es simétrica y definida positiva.
Su representación coordenada cambia con la base, pero el objeto geométrico que
define ---los conjuntos de nivel elipsoidales de $V$--- no cambia.
Así, la cantidad estadística $V$, la forma cuadrática del Capítulo~II y el
tensor de tipo $(0,2)$ del presente capítulo son tres descripciones de la
misma estructura invariante.

\medskip
\hrule
\bigskip

\bigskip
\hrule
\medskip

\textit{Nota: Cubos OLAP y contracción tensorial}

\medskip

En el análisis de datos a gran escala, las observaciones se organizan con
frecuencia a lo largo de múltiples dimensiones categóricas independientes como
tiempo, ubicación, componente del sistema o tipo de medición.
La estructura resultante se llama comúnmente \emph{cubo OLAP} (cubo de
Procesamiento Analítico en Línea, por sus siglas en inglés).

\medskip

Matemáticamente, tal objeto asigna valores numéricos a elementos de un
producto cartesiano
\[
D_1 \times D_2 \times \cdots \times D_n,
\]
donde cada conjunto $D_i$ representa una dimensión de observación.
Tras enumerar estas dimensiones, los datos pueden escribirse como una tabla
multidimensional
\[
F_{i_1 i_2 \dots i_n},
\]
que es precisamente la representación coordenada de un tensor indexado por
$n$ parámetros independientes.

\medskip

En la práctica, solo se observa un pequeño subconjunto de todas las
combinaciones posibles de índices, por lo que la tabla es típicamente dispersa.
Esto refleja de nuevo una distinción enfatizada a lo largo de este capítulo:
las tablas dependen de las coordenadas y las convenciones de almacenamiento
elegidas, mientras que la estructura multilineal subyacente es independiente
de la representación.

\medskip

Las consultas analíticas suelen agregar datos a lo largo de dimensiones
seleccionadas.
Por ejemplo, combinar observaciones ignorando ciertos índices produce cantidades
de la forma
\[
G_{i_1 \dots i_k}
=
\sum_{j_1,\dots,j_m}
F_{i_1 \dots i_k j_1 \dots j_m}.
\]
Desde el punto de vista tensorial, esta operación es una contracción: los
índices correspondientes a las dimensiones agregadas se suman, reduciendo el
orden del tensor.

\medskip

Así, las operaciones conocidas operacionalmente como filtrado, agrupación o
agregación corresponden matemáticamente a proyecciones de un tensor de alta
dimensión sobre subtensores de dimensión inferior.
El cubo OLAP proporciona por tanto una realización concreta de la estructura
tensorial que surge en el análisis de datos moderno.

\medskip

Este ejemplo refuerza un tema recurrente: la contracción tensorial aparece
dondequiera que las relaciones multidimensionales se resumen mediante
agregación sistemática.

\medskip
\hrule
\bigskip
\chapter{Los Logaritmos como Transformaciones que Preservan Estructura}

En los capítulos precedentes, la invariancia bajo cambios de coordenadas fue
identificada como la característica definitoria de los objetos geométricos y
tensoriales.
Un vector, una forma bilineal, un tensor --- ninguno queda definido por sus
componentes numéricas, sino por las relaciones estructurales que preserva bajo
transformación.
El presente capítulo constata que ese mismo principio opera fuera del ámbito
lineal.
El logaritmo es un isomorfismo entre estructuras algebraicas: convierte la
multiplicación en adición preservando íntegramente las relaciones que importan.
El invariante no es una coordenada, una longitud ni una componente tensorial ---
es la propia correspondencia algebraica.

% ------------------------------------------------------------------
%\section*{La Idea Central: una Aplicación que Preserva Estructura}
% ------------------------------------------------------------------

\medskip

En su nivel más elemental, el logaritmo responde a la pregunta:
\emph{¿qué exponente produce un número dado?}
Sin embargo, tanto histórica como matemáticamente,
el poder de los logaritmos radica en algo más profundo:
convierten la multiplicación en adición --- y lo hacen de forma exacta.

\medskip

En el Capítulo~I se definió un isomorfismo como una aplicación entre sistemas
algebraicos que preserva su estructura.
El logaritmo proporciona una realización precisa y global de esa idea.

\medskip

La aplicación
\[
\log:\mathbb{R}^+\to\mathbb{R}
\]
no es una mera biyección; es un isomorfismo entre dos estructuras de grupo:

\begin{enumerate}
    \item Un grupo multiplicativo $(\mathbb{R}^+,\cdot)$.
    \item Un grupo aditivo $(\mathbb{R},+)$.
    \item Un isomorfismo
    \[
    \log:(\mathbb{R}^+,\cdot)\longrightarrow(\mathbb{R},+),
    \]
    cuya inversa es $\exp:(\mathbb{R},+)\to(\mathbb{R}^+,\cdot)$.
\end{enumerate}

\medskip

La propiedad definitoria es:
\[
\log(ab)=\log(a)+\log(b),
\qquad
\log(1)=0.
\]

La multiplicación corresponde a la adición; el elemento neutro multiplicativo
corresponde al elemento neutro aditivo; los inversos multiplicativos corresponden
a los inversos aditivos.
Ninguna relación algebraica se pierde --- solo cambia su apariencia.

\medskip

No se trata de un artificio numérico ni de un atajo computacional.
Es una equivalencia estructural: los dos grupos son, desde el punto de vista
algebraico, el mismo objeto con distinta notación.

% ------------------------------------------------------------------
%\section*{Manifestación Física: la Regla de Cálculo}
% ------------------------------------------------------------------

\medskip

Este isomorfismo cobra existencia física en la regla de cálculo.
Las distancias sobre una escala logarítmica representan valores logarítmicos.
Sumar distancias equivale a multiplicar números.

\medskip

La construcción de escalas logarítmicas consiste en aplicar la transformación
\[
x \mapsto \log_{10}(x),
\]
situando cada número a una distancia proporcional a su logaritmo.
Razones iguales corresponden a distancias iguales.
La identidad estructural $\log(ab)=\log(a)+\log(b)$ se convierte en una
operación mecánica: al deslizar una regla sobre otra, la multiplicación se
realiza mediante una adición de longitudes.

\medskip

Las longitudes no son números, pero obedecen las mismas reglas estructurales.
Al explotar esta correspondencia, la regla de cálculo se convierte en un
computador analógico --- un modelo físico de un isomorfismo algebraico.

% ------------------------------------------------------------------
%\section*{El Logaritmo como Área: Interpretación Geométrica}
% ------------------------------------------------------------------

\medskip

Los logaritmos admiten también una definición puramente geométrica,
independiente de la exponenciación.
Por construcción,
\[
\ln(a)=\int_1^a \frac{1}{x}\,dx,
\]
esto es, el área con signo bajo la curva $y=\frac{1}{x}$ entre $1$ y $a$.

\medskip

Esta definición integral hace transparente la propiedad de isomorfismo.
Para $a,b>0$,
\[
\ln(ab)
=\int_1^{ab}\frac{1}{x}\,dx
=\int_1^{a}\frac{1}{x}\,dx+\int_a^{ab}\frac{1}{x}\,dx.
\]
El cambio de variable $x=at$ en la segunda integral da
\[
\int_a^{ab}\frac{1}{x}\,dx=\int_1^{b}\frac{1}{t}\,dt=\ln(b),
\]
y por tanto $\ln(ab)=\ln(a)+\ln(b)$.
El isomorfismo de grupos es un teorema, no un supuesto.

\medskip

La definición integral también proporciona de inmediato la fórmula fundamental
de derivación:
\[
\frac{d}{da}\ln(a)=\frac{1}{a}.
\]

% ------------------------------------------------------------------
%\section*{Tasa de Cambio y Aproximación Lineal}
% ------------------------------------------------------------------

\medskip

El comportamiento de $\frac{d}{da}\ln(a) = \frac{1}{a}$ conecta el logaritmo
con el marco general de las derivadas y la linealización local.

\medskip

Recordemos que para una función $y = f(x)$, la derivada en un punto $a$ es
\[
f'(a) = \lim_{h\to0}\frac{f(a+h)-f(a)}{h},
\]
la tasa de cambio instantánea y, geométricamente, la pendiente de la recta
tangente.
En un entorno de $a$, toda función diferenciable se comporta aproximadamente
como su recta tangente:
\[
f(x) \approx f(a) + f'(a)(x-a).
\]

\medskip

Esta linealidad local es el reflejo, en el ámbito del análisis, del paso de
los tensores a sus componentes tratado en el Capítulo~II: la estructura global
no lineal se vuelve manejable al aproximarse por su mejor modelo lineal en
cada punto.
La derivada sustituye localmente una aplicación no lineal por una lineal;
el logaritmo sustituye globalmente una estructura multiplicativa por una
aditiva.
Ambos son instancias del mismo principio organizador.

% ------------------------------------------------------------------
%\section*{Cálculo Numérico: Newton--Raphson Aplicado al Logaritmo}
% ------------------------------------------------------------------

\medskip

La derivada del logaritmo también permite su cálculo numérico eficiente
mediante el método de Newton--Raphson.
Calcular $\ln(x)$ equivale a resolver la ecuación
\[
e^y = x,
\]
o, de forma equivalente, encontrar el cero de $f(y) = e^y - x$.

\medskip

La iteración de Newton--Raphson reemplaza la ecuación no lineal por una
sucesión de ecuaciones lineales.
Partiendo de una aproximación inicial $y_n$, se linealiza $f$ en torno a $y_n$
y se resuelve:
\[
0 = f(y_n) + f'(y_n)(y_{n+1}-y_n).
\]
Como $f'(y) = e^y$, se obtiene
\[
y_{n+1} = y_n - \frac{e^{y_n}-x}{e^{y_n}}
         = y_n + \frac{x - e^{y_n}}{e^{y_n}}.
\]

Partiendo de $y_0 = 1{,}5$, las iteraciones sucesivas convergen rápidamente a
$\ln(5) \approx 1{,}60944$.

\medskip

La simplicidad de esta iteración no es casual.
Trabajar en base $e$ garantiza que $\frac{d}{dy}e^y = e^y$,
lo que minimiza la carga algebraica y mejora la estabilidad numérica.
Los logaritmos en otras bases se obtienen mediante la fórmula de cambio de base:
\[
\log_{b}(x) = \frac{\ln(x)}{\ln(b)}.
\]

% ------------------------------------------------------------------
%\section*{Algoritmo Histórico: Briggs y la Acumulación Binaria}
% ------------------------------------------------------------------

\medskip

La misma idea estructural subyace a los algoritmos históricos de construcción
de tablas de logaritmos.
Henry Briggs calculó tablas en base~10 acumulando pequeñas correcciones
multiplicativas; un método estrechamente relacionado, adaptado a la aritmética
binaria, se describe en~\cite{knuth}.

\medskip

El algoritmo construye el número objetivo mediante multiplicaciones sucesivas
de la forma
\[
\frac{10^k}{10^k-1},
\]
acumulando los logaritmos correspondientes cada vez que el producto parcial no
supera el objetivo.
Lo que aparece como cómputo es, en rigor, preservación de estructura:
cada paso multiplicativo corresponde a un incremento aditivo en el lado
logarítmico, ejecutando fielmente el isomorfismo $\log(ab) = \log(a) + \log(b)$.

% ------------------------------------------------------------------
%\section*{La Constante $e$ y sus Representaciones}
% ------------------------------------------------------------------

\medskip

La base natural $e$ no es una elección arbitraria.
Es la única base para la que $\frac{d}{dx}e^x = e^x$,
y la única solución a $\int_1^e \frac{1}{x}\,dx = 1$.
Su carácter aritmético queda revelado por varias construcciones independientes.

\medskip

Una definición clásica proviene del interés compuesto:
\[
e = \lim_{n\to\infty}\left(1+\frac{1}{n}\right)^n.
\]

Otra la proporciona su desarrollo en fracción continua:
\[
e = [2;\,1,2,1,1,4,1,1,6,\dots],
\]
que en forma extendida se escribe
\[
e
=
2 + \cfrac{1}{1 + \cfrac{1}{2 + \cfrac{1}{1 + \cfrac{1}{1 + \cfrac{1}{4 + \cdots}}}}}.
\]

Estas representaciones convergen rápidamente y resultan adecuadas para el
cómputo, pero su significado más profundo es estructural:
$e$ es la unidad natural del isomorfismo entre $(\mathbb{R}^+,\cdot)$
y $(\mathbb{R},+)$.

% ------------------------------------------------------------------
%\section*{Dos Linealizaciones Complementarias}
% ------------------------------------------------------------------

\medskip

El logaritmo transforma el crecimiento no lineal en desplazamiento lineal.
Convierte la multiplicación en adición, los productos en sumas
y las escalas exponenciales en escalas lineales.

\medskip

Se hacen así visibles dos manifestaciones complementarias del principio
invariante.
El logaritmo es un isomorfismo \emph{global}:
establece una equivalencia estructural perfecta entre dos mundos algebraicos,
válida en todo $\mathbb{R}^+$.
La derivada, en cambio, proporciona una linealización \emph{local}:
aproxima una transformación arbitraria por su mejor modelo lineal en cada punto.

\medskip

En ambos casos, la complejidad se hace inteligible al expresarse mediante
relaciones lineales.
Por eso los logaritmos aparecen simultáneamente en el análisis, la geometría,
los métodos numéricos y los instrumentos de cómputo.
No pertenecen a un único dominio.
Preservan estructura a través de todos ellos.

\backmatter
\clearpage
\phantomsection
\addcontentsline{toc}{chapter}{Bibliografía}

\begin{thebibliography}{99}

\bibitem{iribarren}
Iribarren Ignacio L., \textit{Álgebra Lineal}, Editorial Equinoccio, 2009.

\bibitem{halmos}
Halmos, Paul R., \textit{Finite-Dimensional Vector Spaces}, Springer-Verlag, 1987.

\bibitem{iribarren2}
Iribarren, Ignacio L., \textit{Topología de Espacios Métricos}, Editorial Limusa-Wiley, 1973.

\bibitem{cohen}
Cohen, Leon and Ehrlich, Gertrude, \textit{The Structure of the Real Number System}, D. Van Nostrand Company, 1963.

\bibitem{wilder}
Wilder, Raymond L., \textit{Introduction to The Foundations of Mathematics}, Dover Publications, Inc., 2012.

\bibitem{greub}
Greub, Werner, \textit{Linear Algebra}, Springer-Verlag, 1975.

\bibitem{oxford}
Lukas, Andre, \textit{The Oxford Linear Algebra for Scientists}, Oxford University Press, 2022.

\bibitem{apostol}
Apostol, Tom M., \textit{Calculus Volume I}, John Wiley \& Sons Inc., 1967.

\bibitem{knuth}
Knuth, Donald E., \textit{The Art of Computer Programming Volume I: Fundamental Algorithms}, Addison-Wesley, 1997.

\bibitem{knuth2}
Knuth, Donald E., \textit{The Art of Computer Programming Volume II: Seminumerical Algorithms}, Addison-Wesley, 1981.

\bibitem{knuth3}
Knuth, Donald E., \textit{The Art of Computer Programming Volume I, Fascicle 1: MMIX—A RISC Computer for the New Millennium}, Addison-Wesley, 2005.

\bibitem{feynman}
Feynman Richard P., Leighton Robert, Sands Matthew, \textit{The Feynman Lectures on Physics, Volume II}, California Institute of Technology, 2010.

\bibitem{kolecki}
Kolecki, Joseph C., \textit{An Introduction to Tensors for Students of Physics and Engineering}, NASA/TM-2002-211716, 2002.

\bibitem{kolecki2}
Kolecki, Joseph C., \textit{Foundations of Tensor Analysis for Students of Physics and Engineering With an Introduction
to the Theory of Relativity}, NASA/TP—2005-213115, 2005.

\bibitem{kung}
Kung, H. T., Liserson, Charles L., \textit{Introduction to VLSI Systems - Algorithms for VLSI Processor Arrays}, ed. C. Mead and L. Conway, Addison-Wesley, Reading, MA, 1980.

\bibitem{Xu}
Shaofu Xu, Jing Wang, Haowen Shu, Zhike Zhang, Sicheng Yi, Bowen Bai, XingjunWang, Jianguo Liu, Weiwen Zou, \textit{Optical coherent dot-product chip for sophisticated deep learning regression}, Nature - Light: Science \& Applications, 2021.

\end{thebibliography}

\end{document}